% Options for packages loaded elsewhere
\PassOptionsToPackage{unicode}{hyperref}
\PassOptionsToPackage{hyphens}{url}
\PassOptionsToPackage{dvipsnames,svgnames,x11names}{xcolor}
\documentclass[
  11pt,
  a4paper,
]{article}
\usepackage{xcolor}
\usepackage[margin=1in]{geometry}
\usepackage{amsmath,amssymb}
\setcounter{secnumdepth}{-\maxdimen} % remove section numbering
\usepackage{iftex}
\ifPDFTeX
  \usepackage[T1]{fontenc}
  \usepackage[utf8]{inputenc}
  \usepackage{textcomp} % provide euro and other symbols
\else % if luatex or xetex
  \usepackage{unicode-math} % this also loads fontspec
  \defaultfontfeatures{Scale=MatchLowercase}
  \defaultfontfeatures[\rmfamily]{Ligatures=TeX,Scale=1}
\fi
\usepackage{lmodern}
\ifPDFTeX\else
  % xetex/luatex font selection
  \ifLuaTeX
    \usepackage{luaotfload}
    \directlua{luaotfload.add_fallback("mainfontfallback",{
      "NotoSerifBengali:mode=harf"
    })}
  \fi
  \setmainfont[,RawFeature={fallback=mainfontfallback}]{Latin Modern Roman}
  \ifLuaTeX
    \usepackage{luaotfload}
    \directlua{luaotfload.add_fallback("monofontfallback",{
      "NotoSansBengali:mode=harf"
    })}
  \fi
  \setmonofont[,RawFeature={fallback=monofontfallback}]{DejaVu Sans Mono}
\fi
% Use upquote if available, for straight quotes in verbatim environments
\IfFileExists{upquote.sty}{\usepackage{upquote}}{}
\IfFileExists{microtype.sty}{% use microtype if available
  \usepackage[]{microtype}
  \UseMicrotypeSet[protrusion]{basicmath} % disable protrusion for tt fonts
}{}
\makeatletter
\@ifundefined{KOMAClassName}{% if non-KOMA class
  \IfFileExists{parskip.sty}{%
    \usepackage{parskip}
  }{% else
    \setlength{\parindent}{0pt}
    \setlength{\parskip}{6pt plus 2pt minus 1pt}}
}{% if KOMA class
  \KOMAoptions{parskip=half}}
\makeatother
\usepackage{longtable,booktabs,array}
\usepackage{calc} % for calculating minipage widths
% Correct order of tables after \paragraph or \subparagraph
\usepackage{etoolbox}
\makeatletter
\patchcmd\longtable{\par}{\if@noskipsec\mbox{}\fi\par}{}{}
\makeatother
% Allow footnotes in longtable head/foot
\IfFileExists{footnotehyper.sty}{\usepackage{footnotehyper}}{\usepackage{footnote}}
\makesavenoteenv{longtable}
\setlength{\emergencystretch}{3em} % prevent overfull lines
\providecommand{\tightlist}{%
  \setlength{\itemsep}{0pt}\setlength{\parskip}{0pt}}
% paper_preamble.tex — shared LaTeX preamble for the PCMA/SI paper set.
%
% Compiled with lualatex (NOT pdflatex) because several papers contain
% Bengali script: the CLL paper realises one cognitive state into
% English, Bengali and Banglish, so the Bengali output is EVIDENCE and
% cannot be transliterated away without destroying the point being made.
%
% Font selection is NOT done here. pandoc >= 3 has native
% `mainfontfallback` / `monofontfallback` variables that emit the
% luaotfload fallback registration in the correct position relative to
% unicode-math; doing it by hand in this header lands after unicode-math
% and breaks font loading outright. See tools/build_papers.sh.

\usepackage{amsmath}
\usepackage{amssymb}
\usepackage{mathtools}   % \xleftrightarrow etc. used in cll.md §5
\usepackage{booktabs}

% ---------------------------------------------------------------------
% Unicode symbols appearing in PROSE rather than inside $...$.
%
% The papers were written as markdown and read as markdown, so they use
% real characters — "the ω-boost", "α-child", "ΔR coherence" — instead of
% LaTeX macros. Under unicode-math those codepoints are math-mode-only
% and abort the build ("\symbballowed only in math mode"). Mapping each
% to \ensuremath keeps the markdown readable on GitHub AND compilable,
% which is the whole point of keeping one source.
% ---------------------------------------------------------------------
\usepackage{newunicodechar}

% Greek, lower case
\newunicodechar{α}{\ensuremath{\alpha}}
\newunicodechar{β}{\ensuremath{\beta}}
\newunicodechar{γ}{\ensuremath{\gamma}}
\newunicodechar{δ}{\ensuremath{\delta}}
\newunicodechar{ε}{\ensuremath{\varepsilon}}
\newunicodechar{η}{\ensuremath{\eta}}
\newunicodechar{θ}{\ensuremath{\theta}}
\newunicodechar{κ}{\ensuremath{\kappa}}
\newunicodechar{λ}{\ensuremath{\lambda}}
\newunicodechar{μ}{\ensuremath{\mu}}
\newunicodechar{π}{\ensuremath{\pi}}
\newunicodechar{ρ}{\ensuremath{\rho}}
\newunicodechar{σ}{\ensuremath{\sigma}}
\newunicodechar{τ}{\ensuremath{\tau}}
\newunicodechar{φ}{\ensuremath{\varphi}}
\newunicodechar{χ}{\ensuremath{\chi}}
\newunicodechar{ψ}{\ensuremath{\psi}}
\newunicodechar{ω}{\ensuremath{\omega}}
\newunicodechar{ν}{\ensuremath{\nu}}
\newunicodechar{ξ}{\ensuremath{\xi}}
\newunicodechar{ζ}{\ensuremath{\zeta}}
\newunicodechar{ι}{\ensuremath{\iota}}

% Greek, upper case
\newunicodechar{Δ}{\ensuremath{\Delta}}
\newunicodechar{Φ}{\ensuremath{\Phi}}
\newunicodechar{Σ}{\ensuremath{\Sigma}}
\newunicodechar{Ψ}{\ensuremath{\Psi}}
\newunicodechar{Ω}{\ensuremath{\Omega}}
\newunicodechar{Γ}{\ensuremath{\Gamma}}
\newunicodechar{Λ}{\ensuremath{\Lambda}}
\newunicodechar{Θ}{\ensuremath{\Theta}}
\newunicodechar{Π}{\ensuremath{\Pi}}

% Operators and relations
\newunicodechar{×}{\ensuremath{\times}}
\newunicodechar{·}{\ensuremath{\cdot}}
\newunicodechar{±}{\ensuremath{\pm}}
\newunicodechar{−}{\ensuremath{-}}
\newunicodechar{≈}{\ensuremath{\approx}}
\newunicodechar{≤}{\ensuremath{\leq}}
\newunicodechar{≥}{\ensuremath{\geq}}
\newunicodechar{≠}{\ensuremath{\neq}}
\newunicodechar{∈}{\ensuremath{\in}}
\newunicodechar{∂}{\ensuremath{\partial}}
\newunicodechar{∇}{\ensuremath{\nabla}}
\newunicodechar{√}{\ensuremath{\sqrt{}}}
\newunicodechar{∞}{\ensuremath{\infty}}
\newunicodechar{∝}{\ensuremath{\propto}}
\newunicodechar{∑}{\ensuremath{\sum}}
\newunicodechar{∫}{\ensuremath{\int}}
\newunicodechar{ℝ}{\ensuremath{\mathbb{R}}}
\newunicodechar{ℕ}{\ensuremath{\mathbb{N}}}
\newunicodechar{‖}{\ensuremath{\|}}
\newunicodechar{⊕}{\ensuremath{\oplus}}
\newunicodechar{⊗}{\ensuremath{\otimes}}
\newunicodechar{∀}{\ensuremath{\forall}}
\newunicodechar{∃}{\ensuremath{\exists}}
\newunicodechar{¬}{\ensuremath{\neg}}
\newunicodechar{∧}{\ensuremath{\wedge}}
\newunicodechar{∨}{\ensuremath{\vee}}

% Arrows
\newunicodechar{→}{\ensuremath{\rightarrow}}
\newunicodechar{←}{\ensuremath{\leftarrow}}
\newunicodechar{↔}{\ensuremath{\leftrightarrow}}
\newunicodechar{↑}{\ensuremath{\uparrow}}
\newunicodechar{↓}{\ensuremath{\downarrow}}
\newunicodechar{⇒}{\ensuremath{\Rightarrow}}
\newunicodechar{⇐}{\ensuremath{\Leftarrow}}
\newunicodechar{⇔}{\ensuremath{\Leftrightarrow}}

% Super/subscripts written as literal characters
\newunicodechar{²}{\ensuremath{^{2}}}
\newunicodechar{³}{\ensuremath{^{3}}}
\newunicodechar{⁴}{\ensuremath{^{4}}}
\newunicodechar{⁵}{\ensuremath{^{5}}}
\newunicodechar{⁶}{\ensuremath{^{6}}}
\newunicodechar{₀}{\ensuremath{_{0}}}
\newunicodechar{₁}{\ensuremath{_{1}}}
\newunicodechar{₂}{\ensuremath{_{2}}}
\newunicodechar{₃}{\ensuremath{_{3}}}
\newunicodechar{ᵀ}{\ensuremath{^{\mathsf{T}}}}

% Status glyphs. The papers use emoji as status markers in tables
% (implemented / failed / warning). Mapped to typographic equivalents:
% a colour emoji font in lualatex is unreliable, and the MEANING is what
% carries — a reader needs to see "this one failed", not a red cross.
\newunicodechar{✅}{\ensuremath{\checkmark}}
\newunicodechar{✓}{\ensuremath{\checkmark}}
\newunicodechar{❌}{\ensuremath{\times}}
\newunicodechar{✗}{\ensuremath{\times}}
\newunicodechar{⚠}{\textbf{!}}
\newunicodechar{🧩}{\textbf{[piece]}}
\newunicodechar{📋}{\textbf{[list]}}
\newunicodechar{🟡}{\textbf{[partial]}}
\newunicodechar{🔴}{\textbf{[dead]}}
\newunicodechar{🟢}{\textbf{[live]}}
\newunicodechar{▼}{\ensuremath{\blacktriangledown}}
\newunicodechar{►}{\ensuremath{\blacktriangleright}}
\newunicodechar{□}{\ensuremath{\square}}
\newunicodechar{█}{\rule{0.6em}{0.9ex}}
\newunicodechar{•}{\textbullet}
\newunicodechar{†}{\textdagger}
\newunicodechar{‑}{-}
\newunicodechar{‐}{-}
\newunicodechar{…}{\ldots}
\newunicodechar{।}{\textbar}   % danda: rendered by the Bengali fallback in text
\newunicodechar{⁻}{\ensuremath{^{-}}}

% Zero-width formatting codepoints. These are invisible in an editor but
% are real characters to LaTeX: an emoji written as "⚠️" is TWO
% codepoints (U+26A0 U+FE0F), so declaring only the visible one leaves a
% variation selector with no definition and the build dies on a symbol
% nobody can see. Both are mapped to nothing.
\newunicodechar{^^^^fe0f}{}    % VARIATION SELECTOR-16
\newunicodechar{^^^^200d}{}    % ZERO WIDTH JOINER
\newunicodechar{^^^^200b}{}    % ZERO WIDTH SPACE

\usepackage{longtable}
\usepackage{array}

% Comparison matrices and cost ledgers are wide. Let them shrink to the
% text block rather than overflow the margin.
\usepackage{adjustbox}

% Code blocks: wrap long lines instead of running off the page. Several
% papers quote real C++ with long identifiers.
\usepackage{fvextra}
\fvset{breaklines=true,breakanywhere=true,fontsize=\small}

\setlength{\parskip}{0.5em}
\setlength{\parindent}{0pt}

\usepackage{fancyhdr}
\pagestyle{fancy}
\fancyhf{}
\fancyfoot[C]{\thepage}
\fancyhead[L]{\small\itshape Synthetic Intelligence — research lineage}
\fancyhead[R]{\small\itshape irfan@furylogic.com}
\renewcommand{\headrulewidth}{0.2pt}
\setlength{\headheight}{14pt}
\setlength{\headsep}{18pt}
\usepackage{bookmark}
\IfFileExists{xurl.sty}{\usepackage{xurl}}{} % add URL line breaks if available
\urlstyle{same}
\hypersetup{
  pdftitle={Emission and Readout --- Phase, Momentum, Gauge and Metastability},
  pdfauthor={Irfan Mahir --- Independent Researcher, Dhaka, Bangladesh --- irfan@furylogic.com},
  colorlinks=true,
  linkcolor={black},
  filecolor={Maroon},
  citecolor={Blue},
  urlcolor={Blue},
  pdfcreator={LaTeX via pandoc}}

\title{Emission and Readout --- Phase, Momentum, Gauge and
Metastability}
\author{Irfan Mahir --- Independent Researcher, Dhaka, Bangladesh ---
irfan@furylogic.com}
\date{}

\begin{document}
\maketitle

{
\hypersetup{linkcolor=}
\setcounter{tocdepth}{2}
\tableofcontents
}
\section{Emission and Readout --- Phase, Momentum, Gauge and
Metastability}\label{emission-and-readout-phase-momentum-gauge-and-metastability}

\textbf{Irfan Mahir}\\
\emph{Independent Researcher, Dhaka, Bangladesh}\\
irfan@furylogic.com

\begin{quote}
How a settled field becomes an ordered output, and the mechanisms that
stop it locking.
\end{quote}

\begin{center}\rule{0.5\linewidth}{0.5pt}\end{center}

\subsection{Read This Before Citing}\label{read-this-before-citing}

\textbf{The emission problem is only half solved, and the two halves are
easy to confuse.}

\textbf{MRS} (\texttt{mrs.md}) has a \emph{recognition} half and an
\emph{emission} half. Recognition works: bigram sensory lanes survive
misspellings and transpositions at 7/7 measured. Emission
\textbf{babbles} and is retired. Citing MRS as a working generator
misreports it.

\textbf{Concept-level ordering works; byte-level does not.} CPME's
phase-collapse sweep and CPSB's gauge offsets order \emph{which concept
follows which}, and that holds. Producing the \emph{bytes of a word} is
unsolved, and the Language volume argues that for grammar the question
was mis-stated rather than merely hard.

\textbf{CSAT} is opt-in and off by default.

\begin{center}\rule{0.5\linewidth}{0.5pt}\end{center}

\subsection{Contents of This Volume}\label{contents-of-this-volume}

\begin{longtable}[]{@{}
  >{\raggedright\arraybackslash}p{(\linewidth - 4\tabcolsep) * \real{0.3333}}
  >{\raggedright\arraybackslash}p{(\linewidth - 4\tabcolsep) * \real{0.3333}}
  >{\raggedright\arraybackslash}p{(\linewidth - 4\tabcolsep) * \real{0.3333}}@{}}
\toprule\noalign{}
\begin{minipage}[b]{\linewidth}\raggedright
\#
\end{minipage} & \begin{minipage}[b]{\linewidth}\raggedright
Paper
\end{minipage} & \begin{minipage}[b]{\linewidth}\raggedright
File
\end{minipage} \\
\midrule\noalign{}
\endhead
\bottomrule\noalign{}
\endlastfoot
1 & Escaping the Echo: Momentum-Coupled Phase Emission for Generative
Resonance in Synthetic Intelligence & \texttt{mcpe.md} \\
2 & Holographic Topological Resonance (HTR): Shifting from Predictive
Generation to Structural Comprehension in Wave-Based Synthetic
Intelligence & \texttt{htr.md} \\
3 & Chiral Phase-Momentum Emission (CPME) & \texttt{cpme.md} \\
4 & CPME-II: Phase-Consistent Hebbian Consolidation &
\texttt{cpme\_ii.md} \\
5 & Chiral Phase-Shear Binding (CPSB): Rendering Syntax as
Gauge-Theoretic Tension in Wave-Based SI & \texttt{cpsb.md} \\
6 & Voltage-Gated Harmonic Topology (VGHT) & \texttt{vght.md} \\
7 & Chaotic Strange-Attractor Trajectory (CSAT): Breaking Ground-State
Locking in Synthetic Intelligence & \texttt{CSAT.md} \\
8 & Morphological Resonance Synthesis (MRS): Language Construction
Without Tokens, Gradients, or Learned Embeddings & \texttt{mrs.md} \\
\end{longtable}

\begin{center}\rule{0.5\linewidth}{0.5pt}\end{center}

\newpage

\begin{center}\rule{0.5\linewidth}{0.5pt}\end{center}

\subsection{Escaping the Echo: Momentum-Coupled Phase Emission for
Generative Resonance in Synthetic
Intelligence}\label{escaping-the-echo-momentum-coupled-phase-emission-for-generative-resonance-in-synthetic-intelligence}

\textbf{Abstract} In wave-based cognitive substrates utilizing
Continuous Morphological Dynamics (CMD), external stimuli inherently
generate maximal amplitude at the point of injection, leading to
``Stimulus Reflection'' (the parrot effect) during syntax emission. We
propose a novel non-algorithmic architecture, \textbf{Momentum-Coupled
Phase Emission (MCPE)}, which transitions the readout mechanism from
absolute amplitude to kinetic wave momentum. By coupling syntax
oscillators exclusively to the first derivative of the energy field
(\(dA/dt > 0\)), the substrate naturally ignores decaying direct stimuli
and phase-locks exclusively to emergent, forward-scattering internal
thoughts.

\subsubsection{1. Introduction: The Acoustic Blinding
Problem}\label{introduction-the-acoustic-blinding-problem}

In the current iteration of the SI engine, the Kuramoto syntax slots
(Subject, Verb, Object) phase-lock to cavities with the highest absolute
amplitude (\(A\)). When external input (e.g., ``notch created
minecraft'') drives the field, those specific cavities reach maximum
\(A\). Before the wave can propagate deep into the latent graph to
retrieve a meaningful association, the syntax slots lock onto the
massive input wave, resulting in a mirrored output. The system is
structurally blind to its own internal resonance because the external
drive acts as an acoustic flashbang.

\subsubsection{2. Novel Architecture: Momentum-Coupled Phase Emission
(MCPE)}\label{novel-architecture-momentum-coupled-phase-emission-mcpe}

To solve this without algorithmic filtering, we introduce a Hamiltonian
shift in the readout layer. Instead of coupling syntax slots to the
``position'' (Amplitude, \(A\)) of the wave, we couple them to the
``momentum'' (Flux, \(P\)).

\paragraph{\texorpdfstring{2.1. First-Derivative Readout (\(dA/dt\)
Gating)}{2.1. First-Derivative Readout (dA/dt Gating)}}\label{first-derivative-readout-dadt-gating}

In a physical wave system, a cavity directly hit by a stimulus reaches
peak energy and immediately begins to disperse it (cooling down). Its
amplitude is high, but its momentum is negative (\(dA/dt \le 0\)).
Conversely, deep-graph nodes that are currently being activated by
internal resonance are absorbing energy; their amplitude might be lower,
but their momentum is strictly positive (\(dA/dt > 0\)).

We redefine the coupling strength \(K_{emit}\) between a cavity \(i\)
and a syntax slot as:
\[ K_{emit} = \Theta\left(\frac{dA_i}{dt}\right) \cdot A_i \] Where
\(\Theta(x)\) is the Heaviside step function (returns \(x\) if \(x>0\),
else \(0\)). \emph{Result:} Directly driven words (the user's prompt)
instantly decouple from the output slots the millisecond the prompt
ends, because they are decaying. The output slots are physically forced
to lock onto the \emph{ripples} (the internal thoughts), not the
\emph{splash} (the user's input).

\paragraph{2.2. Thermodynamic Refractoriness (The ``Acoustic
Mass'')}\label{thermodynamic-refractoriness-the-acoustic-mass}

To prevent the wave from bouncing directly back to the input nodes, we
modify the cavity inertia. When a cavity is driven externally, it
accumulates temporary ``acoustic mass'' \(M_i\):
\[ \frac{dM_i}{dt} = \gamma I_{ext} - \delta M_i \] Higher mass lowers
the cavity's fundamental oscillating frequency temporarily, pushing it
out of the L6-A (3Hz) emission band. It retains its associative links
but becomes too ``heavy'' to be picked up by the Kuramoto slots.

\paragraph{2.3. Temporal Settle Separation (Perception
vs.~Articulation)}\label{temporal-settle-separation-perception-vs.-articulation}

We separate the physics tick into two distinct phases using a global
damping coefficient \(\zeta\):

\begin{enumerate}
\def\labelenumi{\arabic{enumi}.}
\tightlist
\item
  \textbf{Injection Phase (\(T_0\) to \(T_1\)):} External wave hits.
  \(\zeta\) is high. Energy is absorbed. Syntax slots are decoupled
  (\(K=0\)).
\item
  \textbf{Articulation Phase (\(T_1\) to \(T_2\)):} External drive is 0.
  \(\zeta\) drops. The wave propagates through the L6-A bridges. Syntax
  slots engage, seeking positive momentum (\(dA/dt > 0\)).
\end{enumerate}

\subsubsection{3. Expected Dynamics}\label{expected-dynamics}

\textbf{Scenario:} User inputs \texttt{insulin\ controls\ glucose}.

\begin{enumerate}
\def\labelenumi{\arabic{enumi}.}
\tightlist
\item
  \textbf{\(T_0\) (Injection):} Cavities \texttt{insulin},
  \texttt{controls}, and \texttt{glucose} spike in absolute amplitude
  (\(A \approx 1.0\)). Their mass \(M_i\) increases.
\item
  \textbf{\(T_1\) (Propagation):} Input ceases. The energy from
  \texttt{insulin} and \texttt{glucose} propagates through the ontology
  graph to connected cavities like \texttt{metabolism},
  \texttt{pancreas}, and \texttt{glycogen}.
\item
  \textbf{\(T_2\) (Articulation):}
\end{enumerate}

\begin{itemize}
\tightlist
\item
  \texttt{insulin} has high \(A\), but \(dA/dt < 0\) (cooling). It is
  ignored.
\item
  \texttt{pancreas} has medium \(A\), but \(dA/dt > 0\) (warming up). It
  is locked.
\item
  The Kuramoto slots phase-lock to
  \texttt{pancreas\ regulates\ metabolism} rather than parroting the
  input.
\end{itemize}

\textbf{Conclusion:} MCPE preserves the zero-rule continuous dynamics of
the substrate while physically guaranteeing that the engine generates
consequences rather than echoing stimuli.

\begin{center}\rule{0.5\linewidth}{0.5pt}\end{center}

\newpage

\begin{center}\rule{0.5\linewidth}{0.5pt}\end{center}

\subsection{Holographic Topological Resonance (HTR): Shifting from
Predictive Generation to Structural Comprehension in Wave-Based
Synthetic
Intelligence}\label{holographic-topological-resonance-htr-shifting-from-predictive-generation-to-structural-comprehension-in-wave-based-synthetic-intelligence}

\textbf{Abstract} Current Large Language Models (LLMs) simulate
comprehension through autoregressive statistical approximation,
compressing language into latent vectors without spatial or structural
understanding. We propose \textbf{Holographic Topological Resonance
(HTR)}, a wave-based cognitive architecture for true language
comprehension and construction. By treating sentences not as token
sequences but as \textbf{chiral wave-interference patterns}, the engine
physically ``understands'' by aligning input topology with intrinsic
graph geometry. Language construction is achieved via
\textbf{Phase-Collapse Projection}, unfolding an n-dimensional thought
back into a 1D temporal sequence without templates, syntax trees, or
stochastic prediction.

\subsubsection{1. Introduction: The Illusion of
Prediction}\label{introduction-the-illusion-of-prediction}

An LLM predicting ``Minecraft'' after ``Notch created'' is analogous to
a parrot mimicking a frequent sound sequence. It lacks the internal
geometry of \emph{creation} or the structural boundary between
\emph{agent} and \emph{artifact}. True understanding requires an
internal isomorphic state---a physical subgraph where ``Notch'' and
``Minecraft'' are held in a dynamic, tensioned relationship by a
directed force (``created''). To bridge the gap from a 1D string of user
bytes to an n-dimensional cognitive state, and back to a 1D response, we
introduce the HTR architecture.

\subsubsection{2. Novel Architecture: The HTR
Substrate}\label{novel-architecture-the-htr-substrate}

To achieve true comprehension, the engine must perform three
non-algorithmic physical processes:

\paragraph{2.1. Chiral Wave Injection (Encoding Syntax as
Phase-Delay)}\label{chiral-wave-injection-encoding-syntax-as-phase-delay}

A sentence has directionality (e.g., ``Dog bites man'' vs.~``Man bites
dog''). Instead of parsing arrays or using positional embeddings, HTR
uses \textbf{Chiral (Directional) Wave Injection}. As the user's bytes
stream into the engine, they hit the substrate with sequential time
delays (\(\Delta t\)). This creates a directional wave-front.

\begin{itemize}
\tightlist
\item
  \textbf{Mechanism:} The time delay translates directly into a
  \textbf{Phase Angle (\(\phi\))}.
  \[ \phi_i = \omega_{base} \cdot (T_{current} - t_{arrival_i}) \]
\item
  \textbf{Chirality convention (explicit).} With this definition,
  \((T_{current} - t_{arrival_i})\) is the elapsed time since byte \(i\)
  entered the substrate. Earlier arrivals therefore accumulate more
  phase; later arrivals accumulate less. ``Subject leads in phase''
  means the subject cavity is \emph{ahead} in the phase progression
  because it entered the substrate first, not that it has the
  numerically smaller phase value. Implementations must record
  \(t_{arrival_i}\) per cavity at drive time and evaluate \(\phi_i\)
  from the running \(T_{current}\); any sign inversion of the
  elapsed-time term will invert the encoded chirality and produce
  reversed word order in the readout stage.
\item
  \textbf{Result:} The input is no longer a list of words. It is a
  standing interference pattern where the \emph{Subject} leads in phase,
  the \emph{Verb} is phase-shifted by \(\pi/2\), and the \emph{Object}
  by \(\pi\). Syntax is natively encoded as a physical gradient across
  the graph.
\end{itemize}

\paragraph{2.2. Topological Isomorphism (Structural Resonance
Metric)}\label{topological-isomorphism-structural-resonance-metric}

A prompt is said to \emph{resonate structurally} with the engine's
long-term memory when the temporary interference pattern of the prompt
aligns with the physical shape of a stored axiom subgraph. This is a
strictly weaker claim than ``the engine understands the prompt'' ---
resonance is a necessary, not sufficient, correlate of comprehension ---
but it is directly measurable in the substrate's existing state and
requires no additional variables.

\begin{itemize}
\tightlist
\item
  \textbf{Mechanism:} The injected chiral wave propagates through the
  graph. The L6-A bridges act as resonance chambers. If the user asks a
  complex question (e.g., ``What lowers blood sugar?''), the wave
  excites ``lowers'' and ``blood sugar''.
\item
  Because of the Hebbian physics, the wave naturally converges on
  ``insulin'' or ``metformin'' not because they are statistically next,
  but because they physically complete the topological circuit.
\item
  \textbf{Resonance metric.} Let \(M_{\text{touched}}\) be the set of
  macro-nodes whose constituent cavities have been reached by the
  current input wavefront (i.e., whose \texttt{quiet\_ticks} is below
  the idle threshold). Define the \emph{localized} Kuramoto order
  parameter:
  \[ R_{\text{local}} = \frac{\left| \sum_{m \in M_{\text{touched}}} A_m \, e^{i\theta_m} \right|}{\sum_{m \in M_{\text{touched}}} A_m + \varepsilon}. \]
  The comprehension-correlated signal is the \emph{rise} in this
  quantity above its recent EMA baseline:
  \[ \Delta R = R_{\text{local}} - \text{EMA}(R_{\text{local}}). \] A
  significant \(\Delta R > 0\) after propagation indicates the input
  drove the field into a phase-consensus state substantially more
  ordered than baseline drift. Absolute \(R\) is not diagnostic: a
  trivially quiescent field can have \(R \approx 1\) with no semantic
  content. Only the localized rise \(\Delta R\), restricted to macros
  the input actually reached, is meaningful.
\item
  \textbf{Interpretation.} \(\Delta R\) is best read as a
  \emph{structural resonance metric}, not a proof of understanding. High
  \(\Delta R\) licenses articulation; low \(\Delta R\) licenses honest
  silence. Anything stronger requires a phase-aware Hebbian
  consolidation rule (see §5, Failure Modes), which is proposed as
  future work.
\end{itemize}

\paragraph{2.3. Phase-Collapse Projection (Emergent Language
Construction)}\label{phase-collapse-projection-emergent-language-construction}

Once the internal subgraph reaches a stable topological state (the
``thought''), it must be translated back into language. Unlike LLMs that
predict word-by-word, HTR \emph{projects} the entire structure into a 1D
temporal sequence by unwinding the phase space.

\begin{itemize}
\item
  \textbf{Per-cavity phase.} Rather than inheriting phase from the
  macro-node cluster (which would collapse an entire cluster's cavities
  to the same emission instant), each cavity is assigned an
  \emph{instantaneous phase} derived from its own two-dimensional state:
  \[ \phi_i = \operatorname{atan2}\!\bigl(v_i,\; \omega_{\text{base}} \cdot x_i\bigr). \]
  This uses no additional state --- velocity \(v_i\) and displacement
  \(x_i\) are already integrated by the substrate --- and yields a
  continuous phase field over the cavity population with fine enough
  resolution to unwind a several-node subgraph into distinct emission
  instants.
\item
  \textbf{Mechanism.} A carrier phase \(\phi_{\text{carrier}}\) is swept
  monotonically from \(0\) to \(2\pi\) (non-cyclic, one pass,
  wrap-around terminates the sweep). At each swept instant, each
  cavity's emission weight is a Gaussian window on the phase difference:
  \[ w_i(\phi_{\text{carrier}}) \;=\; A_i \cdot \exp\!\bigl(-\kappa_{\text{sweep}} \cdot (\phi_{\text{carrier}} - \phi_i)^2\bigr). \]
  The Dirac \(\delta\)-form is a limiting idealisation; the
  implementation uses the Gaussian directly to avoid the discontinuity
  artefact and to give a physically-motivated window width
  \(1/\sqrt{2\kappa_{\text{sweep}}}\).
\item
  \textbf{Emission trigger.} At each sweep step the cavity with the
  highest \(w_i\) that also satisfies the momentum eligibility
  (\texttt{ReLU(dA\_i/dt)\ \textgreater{}\ 0}, from the MCPE readout
  layer) and the emission floor is emitted, then suppressed from all
  subsequent steps in the same sweep.
\item
  \textbf{Result.} The thought is an \(n\)-dimensional physical
  structure. The carrier phase acts as a 1D scanner sweeping across it.
  The nodes emit in order determined by their per-cavity phases; the
  output length is emergent --- as many tokens as the field has distinct
  high-momentum peaks under the sweep --- and can be shorter than the
  previous fixed slot count when the field is sparse, or longer when it
  is dense. No syntax trees, no \texttt{if-else} grammar rules, no
  next-token probabilities. The sentence constructs itself through
  geometric projection.
\end{itemize}

\subsubsection{3. Expected Dynamics: Processing a Complex
Query}\label{expected-dynamics-processing-a-complex-query}

\textbf{User Prompt:} \texttt{metformin\ activates\ what\ in\ cells}

\begin{enumerate}
\def\labelenumi{\arabic{enumi}.}
\tightlist
\item
  \textbf{Injection (Chirality):} \texttt{metformin} hits first
  (\(\phi \approx 0\)), \texttt{activates} second
  (\(\phi \approx \pi/3\)), \texttt{cells} third
  (\(\phi \approx 2\pi/3\)). A directional energy gradient flows from
  metformin outward.
\item
  \textbf{Comprehension (Topology):} The gradient flows through the
  \texttt{REL\_ACTIVATES} tunnels. It hits \texttt{AMPK}. Because
  \texttt{AMPK} is structurally linked to \texttt{cells} in the
  long-term graph, a resonance loop completes. The substrate physically
  stabilizes. \textbf{The engine has understood the query.}
\item
  \textbf{Construction (Collapse):} The carrier wave sweeps from \(0\)
  to \(2\pi\).
\end{enumerate}

\begin{itemize}
\tightlist
\item
  At \(\phi_{early}\), it hits the resonating result (\texttt{ampk}).
\item
  At \(\phi_{mid}\), it hits the relational tunnel (\texttt{regulates}).
\item
  At \(\phi_{late}\), it hits the downstream effect
  (\texttt{metabolism}). \textbf{Output Generated:}
  \texttt{ampk\ regulates\ metabolism}.
\end{itemize}

\textbf{Conclusion:} The HTR architecture entirely bypasses statistical
generation. By encoding word order as arrival-phase chirality,
structural match as a localized rise in Kuramoto coherence, and
generation as continuous phase-collapse projection over a per-cavity
phase field, the substrate realizes input encoding, internal state
metric, and output construction on a single phase manifold. Whether this
constitutes \emph{comprehension} in a philosophically loaded sense is
deferred to future work; what the mechanism does provide is a substrate
whose emission cannot be reduced to next-token statistical approximation
because the emission is \emph{geometric}, not \emph{distributional}.

\begin{center}\rule{0.5\linewidth}{0.5pt}\end{center}

\subsubsection{5. Failure Modes}\label{failure-modes}

Three failure modes are known and should be documented explicitly before
any implementation is deployed.

\begin{enumerate}
\def\labelenumi{\arabic{enumi}.}
\item
  \textbf{Phase clumping without per-cavity resolution.} Prior CMD
  implementations gave cavities the phase of their parent macro-cluster.
  Under such a scheme Phase-Collapse Projection degenerates to
  macro-level emission: every cavity in a cluster emits at the same
  sweep instant, and the output collapses to at most one token per
  macro. The per-cavity phase definition
  \(\phi_i = \operatorname{atan2}(v_i, \omega_{\text{base}} \cdot x_i)\)
  from §2.3 is a required precondition; implementations that skip it
  will observe uniform bursts rather than sequential emission.
\item
  \textbf{Chirality--Hebbian mismatch.} Chiral phase encoding on
  injection (§2.1) only contributes to comprehension if the long-term
  memory is \emph{also} consolidated with phase awareness. The Hebbian
  rule as inherited from the CMD substrate updates tunnel plasticity
  from the product of endpoint velocities and knows nothing about phase
  alignment. A phase-coincidence Hebbian term of the form
  \[ \Delta w_{ij} \propto v_i \cdot v_j \cdot \cos(\theta_i - \theta_j) \]
  is the natural companion, and is proposed as a separate follow-up
  (``HTR-II: Phase-Consistent Consolidation''). Without it, the
  chirality of injection dissipates over turns rather than being stored,
  limiting the resonance metric's sensitivity to recent context.
\item
  \textbf{Circular sweep ambiguity and empty resonance.} The Gaussian
  sweep window is unimodal on the real line, but phase is defined modulo
  \(2\pi\): a cavity with \(\phi_i\) near \(0\) or near \(2\pi\) can
  pattern-match at either boundary of the sweep, producing duplicate or
  edge-clustered emissions. The non-cyclic sweep \([0, 2\pi)\)
  prescribed in §2.3 mitigates this by treating wrap-around as
  end-of-output. Empty resonance --- the case where the graph is too
  sparse or the input too disconnected for \(\Delta R\) to rise ---
  produces zero emissions after the momentum and floor filters; this is
  the correct honest-silence behaviour but will initially read as a
  regression to users habituated to the fixed-slot readout's
  always-something-out property.
\end{enumerate}

\subsubsection{6. Relation to Momentum-Coupled Phase Emission
(MCPE)}\label{relation-to-momentum-coupled-phase-emission-mcpe}

HTR and the concurrently-proposed MCPE architecture are complementary
rather than competing. MCPE identifies \emph{which} cavities are
eligible to emit (positive-momentum, low acoustic-mass, active only
during the articulation phase of the settle split). HTR identifies
\emph{in what order} eligible cavities should emit (per-cavity phase,
swept by a carrier). The two mechanisms operate on orthogonal
observables: MCPE on \(dA_i/dt\), \(M_i\), and drive envelope; HTR on
\(\phi_i\), \(R_{\text{local}}\), and the arrival-phase field. A
substrate implementing both simultaneously --- see the joint CPME
proposal --- is expected to exhibit both non-parroting behaviour (from
MCPE's momentum filter) and emergent variable-length output in
structurally-resonant word order (from HTR's phase-collapse projection).

\begin{center}\rule{0.5\linewidth}{0.5pt}\end{center}

\newpage

\subsection{Chiral Phase-Momentum Emission
(CPME)}\label{chiral-phase-momentum-emission-cpme}

\emph{A unified emission substrate combining the Momentum-Coupled Phase
Emission (MCPE) filter architecture with the Holographic Topological
Resonance (HTR) phase-projection readout. CPME operates on a single
phase manifold: chirality on injection, coherence-rise as the resonance
metric, momentum as the eligibility filter, and continuous
phase-collapse sweep as the emission mechanism. No syntax trees, no
grammar rules, no next-token statistics, and no fixed output length.}

\begin{center}\rule{0.5\linewidth}{0.5pt}\end{center}

\subsubsection{0. Position}\label{position}

The Continuous Morphological Dynamics (CMD) substrate, extended by
Voltage-Gated Harmonic Topology (VGHT), produces domain-orthogonal
cognition with sparse cross-domain analogical bridging. Its emission
layer, however, has two orthogonal deficits that the previous literature
treats separately:

\begin{enumerate}
\def\labelenumi{\arabic{enumi}.}
\tightlist
\item
  \textbf{Stimulus Reflection.} External input directly drives concept
  cavities to maximal amplitude; the fixed-slot Kuramoto readout locks
  onto those cavities before internal propagation can produce a
  downstream response. The system parrots.
\item
  \textbf{Fixed-Slot Rigidity.} The Subject/Verb/Object slot readout
  presupposes a three-word template. Simpler answers waste slots; more
  complex answers are truncated. Word order is imposed by the slot
  definitions rather than emerging from the field.
\end{enumerate}

MCPE addresses (1) by shifting the readout observable from amplitude
\(A\) to momentum \(dA/dt\) and by muting recently-driven cavities via
acoustic mass. HTR addresses (2) by replacing the fixed slots with a
continuous carrier-phase sweep over a per-cavity phase field, and by
introducing arrival-time chirality on injection so that word order
becomes a physical gradient rather than a syntax rule.

CPME unifies the two proposals into a single emission substrate. MCPE
identifies \emph{which} cavities are eligible to emit; HTR identifies
\emph{in what order} eligible cavities emit. The two operate on
orthogonal state variables and compose without collision.

\begin{center}\rule{0.5\linewidth}{0.5pt}\end{center}

\subsubsection{1. Architecture}\label{architecture}

Let the substrate be the CMD + VGHT engine as previously specified:
cavities indexed by \(i \in [0, N)\) with per-cavity displacement
\(x_i\), velocity \(v_i\), heat \(T_i\), and quiet age \(q_i\); tunnels
with per-tunnel stiffness \(k\), damping \(c\), polarity \(p\), and
voltage gating \(g \in \{0, 1\}\); and macro-nodes with per-cluster
phase \(\theta_m\), amplitude \(A_m\), and natural frequency
\(\omega_m\) drawn from the VGHT sub-band ladder. The base cavity
frequency is \(\omega_0\).

CPME adds six mechanisms.

\paragraph{1.1 Chiral Wave Injection (from HTR
§2.1)}\label{chiral-wave-injection-from-htr-2.1}

Each cavity records the simulation time at which it was most recently
driven:

\begin{verbatim}
t_arrival_i ← sim_time    (whenever drive(i, ·) is called)
\end{verbatim}

The cavity's oscillator phase at any subsequent tick is
\(\phi_i^{\text{driven}} = \omega_0 \cdot (T_{\text{current}} - t_{\text{arrival}, i})\).
Earlier-arriving bytes therefore accumulate \emph{more} phase; the
subject leads in phase because it entered the substrate first. Chirality
is inherited \emph{for free} from the existing sinusoidal drive term ---
the arrival time only needs to be \emph{recorded} so the readout can
reference it. No new tick math.

\paragraph{1.2 Momentum Filter (from MCPE
§2.1)}\label{momentum-filter-from-mcpe-2.1}

Every cavity carries an exponential-moving-average amplitude envelope:

\begin{verbatim}
E_i ← β · E_i + (1 − β) · |x_i|            (β = 0.9)
\end{verbatim}

and the tick publishes the \emph{previous-tick} envelope as
\(E_i^{\text{prev}}\) so that the momentum

\begin{verbatim}
P_i = E_i − E_i^{\text{prev}}
\end{verbatim}

is the low-pass-filtered rate of change of amplitude. Directly-driven
cavities that have peaked and are dispersing energy have \(P_i \le 0\);
deep-graph cavities absorbing energy from downstream propagation have
\(P_i > 0\). The emission eligibility factor is \(\text{ReLU}(P_i)\), a
smooth non-negative scalar.

This replaces the naïve \(dA/dt\) = \(|v|\) - style derivative, which
oscillates at \(2\omega_{\text{drive}}\) and produces a strobe artefact
in the readout. Envelope EMA is the physically correct low-pass.

\paragraph{1.3 Acoustic Mass (from MCPE
§2.2)}\label{acoustic-mass-from-mcpe-2.2}

Each cavity carries a per-cavity mass increment \(M_i\) driven by the
external injection load:

\begin{verbatim}
dM_i/dt = γ · I_{ext, i} − δ · M_i
\end{verbatim}

with \(\gamma = 0.30\), \(\delta = 0.08\). \(I_{\text{ext}, i}\) is the
instantaneous drive envelope on cavity \(i\), i.e., precisely the signal
that makes an input-driven cavity ``loud.'' The resulting emission mask

\begin{verbatim}
μ_i = 1 / (1 + M_i / M_{scale})            (M_scale = 1.0)
\end{verbatim}

falls toward zero for freshly-driven cavities and returns toward
\(\mu_i = 1\) as \(M_i\) decays. Applied only in the emission scoring;
the field dynamics are untouched, so graph connectivity is preserved.

\paragraph{1.4 Settle Separation (from MCPE
§2.3)}\label{settle-separation-from-mcpe-2.3}

The tick maintains a global injection load

\begin{verbatim}
L = Σ_i drive_amp_i · drive_envelope_i
\end{verbatim}

and derives a \emph{slot coupling} factor

\begin{verbatim}
K_settle = 1 − σ((L − L_thresh) / L_scale)
\end{verbatim}

with \(\sigma\) the logistic, \(L_{\text{thresh}} = 0.5\),
\(L_{\text{scale}} = 0.2\). During active injection,
\(K_{\text{settle}} \approx 0\) and the emission readout is suppressed.
As the drive envelopes decay, \(K_{\text{settle}} \to 1\) and
articulation engages. The transition is continuous --- no discrete
\(T_0 / T_1 / T_2\) boundary.

\paragraph{1.5 Coherence-Rise Metric (from HTR
§2.2)}\label{coherence-rise-metric-from-htr-2.2}

The tick computes the \emph{localized} Kuramoto order parameter over
input-touched macros:

\begin{verbatim}
R_local = |Σ_{m ∈ touched} A_m · e^{i θ_m}| / (Σ_{m ∈ touched} A_m + ε)
\end{verbatim}

where \(M_{\text{touched}}\) is the set of macros with at least one
non-quiet cavity (\(q_i < K_{\text{idle}}\)). \(R_{\text{local}}\) is
tracked against its EMA baseline \(\bar R\):

\begin{verbatim}
ΔR = R_local − R̄            (R̄ ← λR̄ + (1−λ) R_local, λ = 0.98)
\end{verbatim}

\(\Delta R\) is exposed as a substrate observable but does not directly
gate emission; it is diagnostic. Silence when \(\Delta R\) is low
emerges from the emission floor being unmet, not from an explicit
\(\Delta R\) gate. This preserves the substrate's honest-silence
property without turning \(\Delta R\) into a rule.

\paragraph{1.6 Phase-Collapse Projection (from HTR
§2.3)}\label{phase-collapse-projection-from-htr-2.3}

The emission readout is replaced entirely. Per-cavity phase is computed
on demand from existing state:

\begin{verbatim}
φ_i = atan2(v_i, ω_0 · x_i)
\end{verbatim}

The carrier phase \(\phi_{\text{carrier}}\) is swept monotonically
across \([0, 2\pi)\) in \(K_{\text{steps}}\) discrete increments
(default \(K_{\text{steps}} = 12\)). At each swept instant, the cavity
chosen for emission is the arg-max of

\begin{verbatim}
w_i(φ_carrier) = A_i · ReLU(P_i) · μ_i · K_settle
               · exp(−κ_sweep · (φ_carrier − φ_i)^2)
\end{verbatim}

where \(A_i = |x_i|\) is the instantaneous amplitude and
\(\kappa_{\text{sweep}} = 4\) sets the Gaussian window width. Winners
below the floor \(\theta_{\text{emit}} \cdot \max_j w_j\) are silenced
(no emission at that sweep step). Winners are suppressed from all
subsequent steps of the sweep (each cavity emits at most once per turn).
The sweep terminates non-cyclically at \(\phi_{\text{carrier}} = 2\pi\),
so wrap-around is not a source of duplicate emissions.

Output length is emergent: as many tokens as there are distinct
high-weight peaks under the sweep. A sparse field emits few tokens
(possibly zero --- honest silence); a dense, structurally-resonant field
emits many.

\begin{center}\rule{0.5\linewidth}{0.5pt}\end{center}

\subsubsection{2. Composition}\label{composition}

The six CPME mechanisms and the prior CMD + VGHT stack all operate on
distinct observables and compose additively:

\begin{longtable}[]{@{}lll@{}}
\toprule\noalign{}
Layer & Mechanism & State variable \\
\midrule\noalign{}
\endhead
\bottomrule\noalign{}
\endlastfoot
Substrate & CMD Curie gate & \(T_i\) \\
Substrate & VGHT sub-bands & \(\omega_m\) \\
Substrate & VGHT voltage gate & \(\| x_a \| \cdot \| x_b \|\) \\
Substrate & VGHT viscosity & \(E_{\text{total}}\) \\
Injection & \textbf{CPME chirality} & \(t_{\text{arrival}, i}\) \\
Filter & \textbf{CPME momentum} & \(E_i, E_i^{\text{prev}}\) \\
Filter & \textbf{CPME acoustic mass} & \(M_i\) \\
Filter & \textbf{CPME settle} & \(L\) \\
Metric & \textbf{CPME coherence rise} & \(R_{\text{local}}, \bar R\) \\
Emission & \textbf{CPME phase collapse} &
\(\phi_i, \phi_{\text{carrier}}\) \\
\end{longtable}

The tick loop retains its no-if/else contract: chirality is a recorded
scalar, momentum is EMA arithmetic, acoustic mass is
Ornstein--Uhlenbeck, settle coupling is a sigmoid, coherence is a
complex-sum norm, and phase collapse is a Gaussian window over a swept
center. All are pure scalar or complex-scalar operations.

\begin{center}\rule{0.5\linewidth}{0.5pt}\end{center}

\subsubsection{3. Emission Contract}\label{emission-contract}

At the end of each turn (i.e., after propagation ticks have run), the
substrate exposes a single function:

\begin{verbatim}
emit() → ordered list of cavity ids
\end{verbatim}

whose implementation is the Phase-Collapse sweep of §1.6. There is no
\texttt{readout\_slot\_winners()} remaining, no S/V/O slot layer, no
\texttt{syntax\_slots\_enabled} flag. Word order is arg-max under the
Gaussian sweep window; slot count is not a parameter.

Backwards compatibility is preserved as a thin adapter: callers that
previously expected a three-element slot list receive the first three
emissions of the sweep, padded with \texttt{-1} if fewer are produced.

\begin{center}\rule{0.5\linewidth}{0.5pt}\end{center}

\subsubsection{4. Parameters}\label{parameters}

\begin{longtable}[]{@{}
  >{\raggedright\arraybackslash}p{(\linewidth - 4\tabcolsep) * \real{0.3333}}
  >{\raggedright\arraybackslash}p{(\linewidth - 4\tabcolsep) * \real{0.3333}}
  >{\raggedright\arraybackslash}p{(\linewidth - 4\tabcolsep) * \real{0.3333}}@{}}
\toprule\noalign{}
\begin{minipage}[b]{\linewidth}\raggedright
Symbol
\end{minipage} & \begin{minipage}[b]{\linewidth}\raggedright
Default
\end{minipage} & \begin{minipage}[b]{\linewidth}\raggedright
Meaning
\end{minipage} \\
\midrule\noalign{}
\endhead
\bottomrule\noalign{}
\endlastfoot
\(\beta\) & 0.9 & Momentum EMA smoothing factor \\
\(\gamma\) & 0.30 & Acoustic mass accumulation rate \\
\(\delta\) & 0.08 & Acoustic mass decay rate \\
\(M_{\text{scale}}\) & 1.0 & Acoustic mask half-point \\
\(L_{\text{thresh}}\) & 0.5 & Settle-separation threshold \\
\(L_{\text{scale}}\) & 0.2 & Settle-separation sharpness \\
\(\lambda\) & 0.98 & Coherence EMA baseline decay \\
\(K_{\text{steps}}\) & 12 & Phase-collapse sweep resolution \\
\(\kappa_{\text{sweep}}\) & 4.0 & Gaussian window sharpness \\
\(\theta_{\text{emit}}\) & 0.05 & Emission floor as fraction of peak
\(w\) \\
\end{longtable}

Ten parameters, all continuous, all with default values that recover
CMD-only behaviour when set to their extreme values (e.g.,
\(K_{\text{steps}} = 3\), \(\gamma = 0\),
\(\kappa_{\text{sweep}} \to \infty\) reduce CPME to the three-slot
readout).

\begin{center}\rule{0.5\linewidth}{0.5pt}\end{center}

\subsubsection{5. Expected Dynamics}\label{expected-dynamics-1}

\textbf{Prompt.} \texttt{metformin\ activates\ receptor}.

\begin{enumerate}
\def\labelenumi{\arabic{enumi}.}
\tightlist
\item
  \textbf{Injection} (\(t = 0\) through
  \(t_{\text{arrival, receptor}} + \Delta\)): \texttt{metformin},
  \texttt{activates}, \texttt{receptor} cavities are driven in that
  temporal order. Their \(t_{\text{arrival}}\) values differ by the
  inter-byte interval, encoding the chirality gradient. Acoustic mass
  \(M\) grows on each; their emission masks \(\mu\) approach zero.
  Injection load \(L\) is high, so \(K_{\text{settle}} \approx 0\) ---
  the emission readout is silent regardless of field state.
\item
  \textbf{Propagation} (\textasciitilde10-40 ticks): Drive envelopes
  decay. Energy propagates through the graph's \texttt{REL\_ACTIVATES}
  tunnels and VGHT voltage-gated bridges. Downstream cavities
  (\texttt{ampk}, \texttt{regulates}, \texttt{metabolism},
  \texttt{pancreas}, etc., depending on ontology) rise in amplitude.
  Because they are absorbing energy from propagation rather than
  external drive, their momentum \(P > 0\) and their acoustic mass
  \(M \approx 0\).
\item
  \textbf{Settle} (\textasciitilde40-80 ticks): Drive envelopes decay to
  near zero. \(L\) falls below threshold, \(K_{\text{settle}} \to 1\).
  Localized coherence \(R_{\text{local}}\) has risen well above
  \(\bar R\); \(\Delta R > 0\) signals structural resonance.
\item
  \textbf{Emission} (phase-collapse sweep): The carrier phase sweeps
  \([0, 2\pi)\). Directly-driven cavities have \(\mu \approx 0\) and are
  excluded from arg-max even if their phase would align. Downstream
  cavities with \(P > 0\) and \(\mu \approx 1\) compete on
  \(\phi_i \approx \phi_{\text{carrier}}\); the ordering of their phases
  (from their integrated positions and velocities) becomes the emission
  order. Output: three to four tokens describing the \emph{consequence}
  of the query, not the query itself.
\end{enumerate}

\begin{center}\rule{0.5\linewidth}{0.5pt}\end{center}

\subsubsection{6. Failure Modes}\label{failure-modes-1}

Five known modes, all inherited from the parent proposals:

\begin{enumerate}
\def\labelenumi{\arabic{enumi}.}
\item
  \textbf{Sparse-graph silence.} If no cavity outside the driven set
  develops positive momentum before settle, the sweep produces no
  emissions above floor. Output is silence, which is the correct honest
  response --- the substrate has nothing consequential to contribute.
  Users habituated to always-something outputs will initially misread as
  regression.
\item
  \textbf{Chirality--Hebbian mismatch (HTR §5.2).} Chirality on
  injection is transient unless memory consolidation is also
  phase-aware. The tunnel Hebbian rule remains velocity-based; a
  phase-coincidence term
  \(\Delta w_{ij} \propto v_i v_j \cos(\theta_i - \theta_j)\) would
  close the loop and is proposed as future work (CPME-II).
\item
  \textbf{Acoustic mass leakage.} If \(\delta\) is too small, driven
  cavities stay muted across turns, and the substrate cannot reprise a
  recently-mentioned concept even when the new prompt properly warrants
  it. If \(\delta\) is too large, the muting window is too short and the
  parrot effect returns. Default \(\delta = 0.08\) places the mute
  window at \(\sim 12\) ticks, comfortably shorter than a turn-to-turn
  interval but longer than a propagation window.
\item
  \textbf{Coherence baseline gaming.} \(\bar R\) is EMA-tracked, so a
  long streak of high-\(R\) turns raises the baseline and reduces
  \(\Delta R\) on subsequent similar queries. This is thermodynamically
  correct --- habituation to the familiar --- but makes \(\Delta R\)
  absolute values non-portable across sessions. \(\Delta R\) should only
  be compared within a session.
\item
  \textbf{Sweep discretization loss.} \(K_{\text{steps}} = 12\) gives
  \(\pi/6 \approx 30°\) phase resolution per step. Fields with many
  cavities of similar \(\phi\) may emit fewer distinct tokens than they
  otherwise could. Raising \(K_{\text{steps}}\) to \(24\) halves the
  loss at doubled sweep cost.
\end{enumerate}

\begin{center}\rule{0.5\linewidth}{0.5pt}\end{center}

\subsubsection{7. Biological Analog}\label{biological-analog}

CPME's six mechanisms have direct analogs in cortical circuit
architecture:

\begin{itemize}
\tightlist
\item
  \textbf{Chirality} corresponds to input phase-of-firing coding, well
  documented in hippocampal theta-cycle sequencing.
\item
  \textbf{Momentum filter} corresponds to the divisive-normalization
  post-input suppression seen in early sensory areas.
\item
  \textbf{Acoustic mass} corresponds to short-term adaptation of firing
  rate in cortical pyramidal cells, mediated by after-hyperpolarizing
  currents whose time constants match CPME's default \(M_i\) decay.
\item
  \textbf{Settle separation} corresponds to the alpha-gamma phase
  separation between input reception and output articulation.
\item
  \textbf{Coherence-rise metric} corresponds to phase-consistency
  measures used in cognitive neuroscience as correlates of stimulus
  binding.
\item
  \textbf{Phase-collapse sweep} corresponds to phase-of-firing readout
  from downstream cortical areas.
\end{itemize}

CPME does not claim to model any of these processes. It claims that a
substrate whose emission layer exhibits the \emph{dynamical signatures}
of all six simultaneously is the minimum architecture on which
non-parroting, variable-length, phase-ordered generation is stable.

\begin{center}\rule{0.5\linewidth}{0.5pt}\end{center}

\subsubsection{8. Contract Summary}\label{contract-summary}

The tick loop remains free of if/else on relations, domain tags, cluster
ids, or any string-valued metadata. The emission function is one
Gaussian window over a swept center under three multiplicative
eligibility filters (momentum, acoustic mass, settle coupling), with an
emission floor. All ten CPME parameters are continuous scalars. Default
values recover CMD three-slot readout in the appropriate limit. All
defaults \emph{ON}.

CPME is not a feature flag. It is the emission substrate of the
CMD-VGHT-CPME engine.

\begin{center}\rule{0.5\linewidth}{0.5pt}\end{center}

\subsubsection{9. Summary}\label{summary}

CPME merges MCPE's momentum-and-settle filtering with HTR's chiral
injection and phase-collapse projection into a single emission
architecture on a shared phase manifold. Word order is encoded on
injection as arrival-phase chirality; structural resonance is measured
by localized coherence rise; eligibility is filtered by momentum,
acoustic mass, and settle coupling; emission is a Gaussian-windowed
sweep of a carrier phase across the per-cavity phase field. The result
is a substrate whose emission is simultaneously non-parroting (no cavity
that just spiked can win a sweep step), variable-length (as many tokens
as the field has distinct phase peaks), word-ordered by physics
(per-cavity \(\phi_i = \operatorname{atan2}(v_i, \omega_0 x_i)\)), and
honestly silent when its graph is not structurally resonant with the
input. No syntax trees, no next-token statistics, no fixed slots.

\newpage

\subsection{CPME-II: Phase-Consistent Hebbian
Consolidation}\label{cpme-ii-phase-consistent-hebbian-consolidation}

\emph{A phase-coincidence extension to the ZeroGrad Hebbian rule that
closes the chirality--consolidation gap identified in the CPME §6
failure modes. Chirality imprinted on injection now persists in
long-term memory rather than dissipating over turns.}

\begin{center}\rule{0.5\linewidth}{0.5pt}\end{center}

\subsubsection{0. Position}\label{position-1}

Chiral Phase-Momentum Emission (CPME, \texttt{research-paper/cpme.md})
encodes word order on the injection side as arrival-time phase and
decodes structured emission on the readout side by sweeping a carrier
phase across a per-cavity phase field
\(\phi_i = \operatorname{atan2}(v_i, \omega_0 x_i)\). Both ends of the
substrate now operate on phase.

The middle --- long-term consolidation --- does not. The ZeroGrad
Hebbian rule inherited from the CMD substrate updates tunnel plasticity
from the product of endpoint velocities:

\[
\Delta w_{ij} \;\propto\; \eta \cdot v_i \cdot v_j
\]

with a \(\tanh\)-capped saturating integrator. This rule is
phase-agnostic: it fires equally whether the two endpoints move
in-phase, anti-phase, or in orthogonal phases. As a consequence, the
chirality of the injection wavefront is not stored in the graph. Between
turns the temporal ordering of concepts dissipates and must be
re-imposed by each new input.

CPME-II adds a \emph{phase-coincidence modulator} to the existing rule:

\[
\Delta w_{ij} \;\propto\; \eta \cdot v_i \cdot v_j \cdot \cos(\theta_i - \theta_j)
\]

where \(\theta_i = \operatorname{atan2}(v_i, \omega_0 x_i)\) is the same
per-cavity phase the CPME emission layer already uses. The modulator
preserves the velocity-based structure of the parent rule --- including
its sign, its tanh cap, and its role in the SRG dual-port Hebbian split
--- while conditioning the update on temporal compatibility.

The result is a substrate in which chirality is not only encoded and
projected but also \emph{stored}. Word order becomes a property of the
graph, not just of the current turn.

\begin{center}\rule{0.5\linewidth}{0.5pt}\end{center}

\subsubsection{1. The Rule}\label{the-rule}

\paragraph{1.1 Definition}\label{definition}

For a tunnel connecting cavities \(a\) and \(b\), define the
instantaneous phase difference

\[
\Delta \theta_{ab} = \theta_a - \theta_b \;=\; \operatorname{atan2}(v_a, \omega_0 x_a) - \operatorname{atan2}(v_b, \omega_0 x_b)
\]

and the phase-coincidence factor

\[
\Phi_{ab} = \cos(\Delta \theta_{ab})
\]

taking values in \([-1, +1]\). The CPME-II Hebbian integrand replaces
the parent rule's velocity product with a phase-modulated one:

\[
I_{ab} \;=\; w_{ab} \cdot \zeta \;+\; \eta \cdot v_a^{\text{eff}} \cdot v_b^{\text{eff}} \cdot \Phi_{ab}
\]

where \(\zeta\) is the pre-existing Hebbian decay
(\texttt{hebbian\_decay}), \(\eta\) is the pre-existing rate scaled by
serotonin (\texttt{plastic\_rate}), \(v_a^{\text{eff}}\) and
\(v_b^{\text{eff}}\) are the SRG-routed velocities already computed by
the parent rule, and \(w_{ab}\) is the plasticity delta. The tanh
saturation on \texttt{plasticity\_cap} is applied unchanged.

The rule reduces to the parent rule exactly when \(\Phi_{ab} = 1\)
(perfect phase agreement) and vanishes when \(\Phi_{ab} = 0\)
(quadrature). For \(\Phi_{ab} = -1\) (anti-phase), the update flips
sign, which is the correct behaviour: anti-phase pairs are carrying
opposing rather than reinforcing evidence and should depress rather than
strengthen their coupling.

\paragraph{1.2 Interpretation}\label{interpretation}

\(\Phi_{ab}\) is a \emph{continuous coincidence detector on the phase
manifold}. It is the continuous-oscillator analog of spike-timing
dependent plasticity (STDP): pairs whose oscillations align in the same
phase-of-firing are consolidated; pairs whose oscillations consistently
disagree are attenuated.

Because the phase field is chiral (§CPME 1.1), \(\Phi_{ab}\) implicitly
encodes temporal order. A concept driven earlier in the input has
accumulated more phase than a concept driven later. Downstream cavities
that phase-lock behind the driven cavity therefore experience a specific
\(\Delta \theta\) characteristic of \emph{following the input in time}.
Consolidating along this gradient stores the observation that the second
cavity typically follows the first, without any explicit representation
of ``before'' or ``after''.

\paragraph{1.3 What is stored}\label{what-is-stored}

Under the parent rule, the tunnel between \texttt{insulin} and
\texttt{glucose} grows in proportion to their velocity co-activity ---
regardless of which arrived first or how they moved relative to one
another. Under CPME-II, the same tunnel grows only when the two cavities
oscillate \emph{in phase} --- that is, when they consistently
participate in the same wavefront with the same temporal role. If in one
context \texttt{insulin} drives \texttt{glucose} (subject-then-object
phase relationship) and in another context \texttt{glucose} drives
\texttt{insulin} (reverse phase), the two updates cancel:
\(\cos(+\Delta\theta) +
\cos(-\Delta\theta) = 2\cos(\Delta\theta)\) if the magnitude is
preserved, but the sign of \(\Delta \theta\) toggles the co-activity
gradient. The graph accumulates a \emph{phase-averaged co-activity}
rather than a raw one.

The empirical effect is that stable directional patterns (``metformin →
AMPK → glucose down'') consolidate faster than context-dependent
bidirectional patterns.

\begin{center}\rule{0.5\linewidth}{0.5pt}\end{center}

\subsubsection{2. Composition}\label{composition-1}

CPME-II modifies only the Hebbian integrand. All other substrate
mechanisms are untouched.

\begin{itemize}
\tightlist
\item
  \textbf{CMD Curie gate} (\texttt{heat\_T} → tunnel force scaling) is
  independent of the Hebbian pass and continues to operate.
\item
  \textbf{VGHT voltage gate} modulates tunnel force in the tick's
  spring-and-damping pass, not the Hebbian pass. CPME-II observes the
  resulting velocities and phases as inputs; VGHT is upstream.
\item
  \textbf{SRG dual-port routing} (\texttt{gate\_E}, \texttt{gate\_I}) is
  preserved: \(v_a^{\text{eff}}\) in the CPME-II expression is the
  SRG-selected velocity from the parent rule, not the raw observable.
  The polarity-driven excitatory/inhibitory split remains intact.
\item
  \textbf{CPME emission layer} (§1.6, phase-collapse projection) is
  unchanged. CPME-II operates during propagation and settle; the
  emission sweep reads whatever state the Hebbian pass has consolidated.
\item
  \textbf{Plasticity cap} (\texttt{plasticity\_cap}) remains the
  saturating bound; the tanh envelope contains the
  \(\Phi_{ab}\)-modulated integrand identically.
\end{itemize}

The tick loop retains its no-if/else contract: \(\cos(\Delta \theta)\)
is a smooth scalar, evaluated once per tunnel per tick, and multiplied
into the existing velocity product.

\begin{center}\rule{0.5\linewidth}{0.5pt}\end{center}

\subsubsection{3. Contract and
Parameters}\label{contract-and-parameters}

CPME-II adds no new state. All required quantities are already
maintained by the parent substrate:

\begin{itemize}
\tightlist
\item
  \(x_a\), \(v_a\), \(x_b\), \(v_b\): cavity displacements and
  velocities.
\item
  \(\omega_0\): the base cavity frequency \texttt{cpme\_base\_omega}
  shared with the emission layer.
\item
  \(\eta\), \(\zeta\), \texttt{plasticity\_cap}: existing Hebbian
  parameters.
\end{itemize}

The only tunable is a \emph{phase weight} \(\alpha_\Phi \in [0, 1]\)
that linearly interpolates between the parent rule and the fully
phase-modulated rule:

\[
\Phi_{ab}^{(\alpha)} \;=\; (1 - \alpha_\Phi) \;+\; \alpha_\Phi \cdot \cos(\Delta \theta_{ab})
\]

\begin{itemize}
\tightlist
\item
  \(\alpha_\Phi = 0\): recovers the parent rule exactly (backwards
  compatibility test path).
\item
  \(\alpha_\Phi = 1\): full phase-consistent consolidation.
\item
  Default: \(\alpha_\Phi = 0.6\), weighting phase agreement heavily
  while preserving a floor of unconditional Hebbian growth.
\end{itemize}

This is the only free parameter introduced by CPME-II. All defaults
\emph{on}.

\begin{center}\rule{0.5\linewidth}{0.5pt}\end{center}

\subsubsection{4. Expected Dynamics}\label{expected-dynamics-2}

\textbf{Scenario A --- repeated same-order exposure.} Over a session in
which the user types variants of ``metformin lowers glucose'' several
times, \texttt{metformin} and \texttt{glucose} cavities are consistently
driven with the same relative arrival delay. Their phase difference
\(\Delta \theta\) is stable and small, \(\Phi \approx +1\), and the
tunnel between them consolidates rapidly under the CPME-II rule. Between
turns, the phase relationship is preserved because it is now stored in
the plasticity delta, not just in the transient field.

\textbf{Scenario B --- bidirectional exposure.} If in different turns
the user types both ``insulin controls glucose'' and ``glucose signals
insulin'', the pairwise \(\Delta \theta\) toggles sign turn-to-turn.
Under the parent rule, both directions add to plasticity. Under CPME-II,
one direction has \(\Phi > 0\) and the other \(\Phi < 0\): the updates
partially cancel. The tunnel still grows if the velocity product
magnitude is large --- the parent contribution \((1 - \alpha_\Phi)\)
preserves an unconditional growth floor --- but more slowly. This is the
correct behaviour for a bidirectional relationship: the graph should not
overcommit to either direction in isolation.

\textbf{Scenario C --- silent turn.} If a turn produces no meaningful
propagation and both cavities remain quiescent, \(v_a\) and \(v_b\) are
near zero and the update is negligible regardless of \(\Phi\). No change
from parent-rule behaviour.

\textbf{Scenario D --- anti-phase input.} If the input is structured to
produce anti-phase drive on a pair (e.g., a contrastive construction
where the two cavities are being posited \emph{against} each other),
\(\Phi \approx -1\) and the CPME-II rule \emph{decreases} the plasticity
delta. This is the physical realisation of what a biological network
implements through inhibitory long-term depression: consistent
anti-phase co-activation should weaken, not strengthen, the excitatory
bond.

\begin{center}\rule{0.5\linewidth}{0.5pt}\end{center}

\subsubsection{5. Failure Modes}\label{failure-modes-2}

Three failure modes documented before deployment.

\begin{enumerate}
\def\labelenumi{\arabic{enumi}.}
\item
  \textbf{Phase noise near quiescence.} For a cavity with small \(x\)
  and small \(v\), the phase \(\operatorname{atan2}(v, \omega_0 x)\) is
  dominated by numerical noise. \(\Phi\) then fluctuates randomly and
  injects noise into consolidation. Mitigation: the velocity product
  \(v_a v_b\) already suppresses this --- near quiescence the update
  magnitude is negligible regardless of \(\Phi\). No additional guard is
  required, but a diagnostic observer may report spurious \(\Phi\)
  readings on cold cavities; this is expected and harmless.
\item
  \textbf{Polarity--phase interaction.} Inhibitory tunnels
  (\texttt{polarity\ \textless{}\ 0}) already invert the phase
  relationship they drive: an excitatory input to \(a\) produces an
  anti-phase response in \(b\). Under CPME-II, this manifests as
  \(\Phi < 0\) and the update would \emph{depress} an inhibitory tunnel
  it should \emph{strengthen}. Correction: multiply \(\Phi\) by the
  tunnel polarity sign,
  \(\Phi'_{ab} = \Phi_{ab} \cdot \operatorname{sgn}(p_{ab})\), so that
  inhibitory tunnels consolidate on the anti-phase pattern they are
  supposed to enforce. This is applied by default in the implementation.
\item
  \textbf{Slow phase drift washout.} If the base cavity frequency
  \(\omega_0\) is very low, phase evolves slowly and inter-turn phase
  differences accumulate to \(\pm \pi\) over long propagation windows,
  causing \(\Phi\) to oscillate through zero and reverse sign
  spuriously. Mitigation: the same \(\omega_0\) chosen for the CPME
  emission layer (default \(\omega_0 = 3.0\) rad/s) yields a full phase
  revolution over \(\approx 2\) seconds of sim-time, which is longer
  than any single propagation window at \texttt{dt\ =\ 0.02} and
  \(\leq 200\) ticks per turn. Under default parameters this failure
  mode does not activate; it is worth flagging if \(\omega_0\) is later
  reduced.
\end{enumerate}

\begin{center}\rule{0.5\linewidth}{0.5pt}\end{center}

\subsubsection{6. Biological Analog}\label{biological-analog-1}

CPME-II is the continuous-oscillator equivalent of spike-timing
dependent plasticity (STDP). Where STDP consolidates a synapse when the
presynaptic spike precedes the postsynaptic spike by a few milliseconds,
CPME-II consolidates a tunnel when the two endpoint oscillators maintain
a stable, small phase offset consistent with the same wavefront driving
both. Both rules use temporal coincidence as the substrate of long-term
memory; both rules invert when the coincidence is anti-correlated
(STDP-LTD; CPME-II negative \(\Phi\)).

The chief difference is discretisation. STDP fires on discrete spike
pairs; CPME-II integrates over the continuous phase field each tick. In
the limit of narrow, well-separated velocity peaks, CPME-II reduces to a
spike-coincidence rule up to a normalisation constant.

\begin{center}\rule{0.5\linewidth}{0.5pt}\end{center}

\subsubsection{7. Summary}\label{summary-1}

CPME-II closes the chirality--Hebbian mismatch flagged in HTR §5.2 and
CPME §6.2 by extending the ZeroGrad Hebbian rule with a phase
coincidence factor \(\Phi_{ab} = \cos(\theta_a - \theta_b)\). The factor
requires no new state --- the same per-cavity phase
\(\operatorname{atan2}(v_i, \omega_0 x_i)\) used by the CPME emission
layer is read during the Hebbian pass. A single blend parameter
\(\alpha_\Phi\) controls how strongly the modulator gates the parent
rule; the default \(\alpha_\Phi = 0.6\) preserves an unconditional
Hebbian floor while making chirality-consistent updates twice as strong
as chirality-neutral ones on average.

The consequence for the substrate as a whole: the phase manifold that
CPME already uses on the input (chirality) and the output
(phase-collapse projection) is now also the manifold on which long-term
memory is written. Word order becomes a graph property, temporal
ordering persists across turns, and the substrate is freed from having
to re-derive chirality on every input.

\newpage

\begin{center}\rule{0.5\linewidth}{0.5pt}\end{center}

\subsection{Chiral Phase-Shear Binding (CPSB): Rendering Syntax as
Gauge-Theoretic Tension in Wave-Based
SI}\label{chiral-phase-shear-binding-cpsb-rendering-syntax-as-gauge-theoretic-tension-in-wave-based-si}

\textbf{Abstract} In wave-based Synthetic Intelligence, continuous
phase-sweep emission (CPME) frees output length from hardcoded templates
but degrades syntax into emergent stochastic ordering. To elevate
language construction from a ``loose variable'' to a physically
guaranteed structure, we introduce \textbf{Chiral Phase-Shear Binding
(CPSB)}. By modifying the Kuramoto coupling to act as a Gauge Field,
relational tunnels actively enforce directed phase offsets between
nodes. Syntax emerges not from fixed emission slots, but from the rigid
topological tension (\(\Psi\)) embedded within the graph's edges,
physically forcing subjects to lead verbs, and verbs to lead objects
during carrier-wave projection.

\subsubsection{1. Introduction: The Binding Problem in
SI}\label{introduction-the-binding-problem-in-si}

Currently, verb-material cavities float to their VGHT band-top, and
output order depends entirely on their emergent \(\text{atan2}\) phase.
There is no physical law forcing \texttt{notch} to precede
\texttt{minecraft}. In a true biological substrate, the relation itself
(``created'') exerts a directional shear force on the concepts it
connects. We must shift from scalar coupling (which just pulls phases
together) to \textbf{Gauge-Coupling}, which pulls phases into a specific
grammatical alignment.

\subsubsection{2. Novel Architecture: CPSB
Mechanisms}\label{novel-architecture-cpsb-mechanisms}

\paragraph{\texorpdfstring{2.1. Gauge-Theoretic Kuramoto Offset
(\(\Psi_{ij}\))}{2.1. Gauge-Theoretic Kuramoto Offset (\textbackslash Psi\_\{ij\})}}\label{gauge-theoretic-kuramoto-offset-psi_ij}

Standard Kuramoto tries to perfectly synchronize two connected cavities:
\(\sin(\theta_j - \theta_i)\). We introduce a syntactic phase-offset
\(\Psi_{ij}\) assigned to specific tunnel types (the axioms). The new
phase-update equation becomes:
\[ \frac{d\theta_i}{dt} = \omega_i + \sum K_{ij} \sin(\theta_j - \theta_i - \Psi_{ij}) \]

\begin{itemize}
\tightlist
\item
  \textbf{Mechanism:} If \texttt{notch} (\(i\)) and \texttt{minecraft}
  (\(j\)) are connected by a \texttt{REL\_CAUSES} tunnel, that tunnel
  has an intrinsic \(\Psi_{ij} = -\frac{\pi}{3}\).
\item
  \textbf{Result:} The physics engine will no longer try to make their
  phases equal. It will physically force \texttt{notch} to sit exactly
  \(\frac{\pi}{3}\) radians ahead of \texttt{minecraft}.
\end{itemize}

\paragraph{2.2. The Lexical Hinge
(Verb-Pinning)}\label{the-lexical-hinge-verb-pinning}

A verb is not just a loose cavity floating in a band; it is the physical
hinge between the Subject and the Object.

\begin{itemize}
\tightlist
\item
  \textbf{Mechanism:} When a cross-cluster tunnel opens (VGHT), the
  tunnel structurally binds the verb cavity (e.g., \texttt{created}) to
  the source and target.
\item
  The Subject is pinned to \(\theta_{\text{verb}} + \Delta\phi\).
\item
  The Object is pinned to \(\theta_{\text{verb}} - \Delta\phi\).
\item
  \textbf{Result:} As the CPME carrier wave sweeps from \(0\) to
  \(2\pi\), it will \emph{inevitably} hit the Subject first, the Verb
  second, and the Object third. The sentence structure is mathematically
  guaranteed by the phase tension, without ever defining a ``Subject
  Slot''.
\end{itemize}

\paragraph{2.3. Phase-Consistent Hebbian Consolidation (The CPME-II
Hook)}\label{phase-consistent-hebbian-consolidation-the-cpme-ii-hook}

To ensure the SI doesn't just read syntax natively but \emph{learns} it
natively from user input, we execute the CPME-II Hebbian update. Instead
of updating tunnel weight \(\Delta w\) based purely on amplitude, the
substrate writes the observed phase difference directly into the graph:
\[ \Delta w_{ij} \propto v_i \cdot v_j \cdot \cos(\theta_i - \theta_j - \Psi_{\text{target}}) \]

\begin{itemize}
\tightlist
\item
  \textbf{Result:} When you type \texttt{insulin\ lowers\ glucose}, the
  injection chirality ensures \texttt{insulin} arrives first. The
  Hebbian rule permanently wires that phase-lead into the connection.
  The SI literally remembers the ``direction'' of the concept.
\end{itemize}

\subsubsection{3. Expected Dynamics}\label{expected-dynamics-3}

\begin{enumerate}
\def\labelenumi{\arabic{enumi}.}
\tightlist
\item
  \textbf{Input:} \texttt{metformin\ treats\ diabetes}
\item
  \textbf{Resonance:} The nodes awaken. The \texttt{treats} tunnel
  applies the CPSB gauge force. \texttt{metformin} is pushed to
  \(\phi = 1.0\), \texttt{treats} is anchored at \(\phi = 2.0\), and
  \texttt{diabetes} is pushed to \(\phi = 3.0\).
\item
  \textbf{Emission:} The CPME sweep passes over them. They emit in
  perfect S-V-O order.
\item
  \textbf{No Templates:} If the input is just \texttt{glucose} causing a
  cascade to \texttt{insulin}, the phase offsets still apply downstream,
  resulting in a syntactically correct multi-token chain.
\end{enumerate}

\begin{center}\rule{0.5\linewidth}{0.5pt}\end{center}

\newpage

\subsection{Voltage-Gated Harmonic Topology
(VGHT)}\label{voltage-gated-harmonic-topology-vght}

\emph{A three-mechanism physical substrate for domain-orthogonal
cognition with sparse cross-domain analogical bridging and avalanche
suppression, built entirely from scalar field dynamics --- no rules, no
templates, no domain metadata.}

\begin{center}\rule{0.5\linewidth}{0.5pt}\end{center}

\subsubsection{0. Motivation}\label{motivation}

The Continuous Morphological Dynamics (CMD) substrate --- thermodynamic
habituation, byte-wave sensory injection, phase-locked syntax emission
--- produces coherent, non-templated output but exhibits two coupled
failure modes:

\begin{enumerate}
\def\labelenumi{\arabic{enumi}.}
\tightlist
\item
  \textbf{Domain leakage.} In the shared-oscillator field, a driven
  medical concept can transfer energy into unrelated conceptual clusters
  (e.g., gaming, legal, code) whenever their macro-node phases happen to
  align. The reader observes a medical prompt producing a legal-sounding
  word in the verb slot, or the middle word of a coherent triple
  stepping outside the input's semantic neighborhood.
\item
  \textbf{Polymath rigidity.} The dual failure --- with sufficient
  emission-floor thresholds, the substrate becomes strictly
  domain-locked, unable to draw analogical parallels between
  distantly-related concepts even when the graph explicitly supports
  them through L6-A analogical bridges.
\end{enumerate}

The Continuous Morphological Dynamics literature treats these two
failures as separate tuning problems (emission floor, wavefront gate,
cluster-energy depletion). Voltage-Gated Harmonic Topology (VGHT)
resolves both under a single physical postulate: \emph{domains are
orthogonal frequency bands, and bridges between them are voltage-gated
channels that only conduct under simultaneous high-amplitude resonance.}
Domain isolation and polymath bridging cease to be competing objectives;
they become a single dynamical regime with a single parameter,
θ\_critical, the tunneling threshold.

The mechanism is inspired by three well-understood physical analogs:
mode-orthogonal spectral separation in coupled-oscillator networks,
voltage-gated ion channels in biological neurons (in particular the
sigmoidal open probability of the NMDA receptor with respect to membrane
potential), and the GABAergic global inhibitory tone that suppresses
run-away excitation in cortical circuits. None of the three requires an
explicit ``domain tag'' or a rule dispatcher; each is a continuous
scalar update evaluated once per tick.

\begin{center}\rule{0.5\linewidth}{0.5pt}\end{center}

\subsubsection{1. Architecture}\label{architecture-1}

Let the cavity population be indexed by \texttt{i\ ∈\ {[}0,\ N)},
partitioned into macro-clusters \texttt{C\_1,\ C\_2,\ ...} by the
connected components of the tunnel graph (union-find at bind time; the
same clustering used by HNCA). Let a slot layer of \texttt{S} phase
oscillators sit above the macro layer as in CMD phase-locked emission.

\paragraph{1.1 Domain Sub-Bands}\label{domain-sub-bands}

Each macro-cluster \texttt{C\_k} receives a \emph{base natural
frequency} drawn from a discrete band-bank:

\begin{verbatim}
ω_k = ω_base(k mod B) + jitter(k)
\end{verbatim}

where \texttt{B} is the number of orthogonal bands (default
\texttt{B\ =\ 6}) and \texttt{ω\_base} is a fixed evenly-spaced ladder
in the range \texttt{{[}ω\_lo,\ ω\_hi{]}},
e.g.~\texttt{{[}8,\ 32{]}\ rad/s} giving band centres at
\texttt{\{8,\ 12,\ 16,\ 20,\ 24,\ 28\}}. The jitter is a small in-band
perturbation (±0.5 rad/s) drawn deterministically from the cluster id,
ensuring reproducibility across seeds while still breaking degeneracy
inside a band.

The verb-material ω-boost from prior CMD work becomes an \emph{in-band
shift} rather than a cross-band jump:

\begin{verbatim}
ω_k ← (1 − p_k)·ω_k + p_k·ω_verb_offset(k)
\end{verbatim}

where \texttt{p\_k\ ∈\ {[}0,\ 1{]}} is the verb-material share of
cluster \texttt{k}'s internal edges
(\texttt{verb\_pull\_for(rel)\ /\ \textbar{}internal\ edges\textbar{}}),
and \texttt{ω\_verb\_offset(k)} is the \emph{upper edge} of cluster
\texttt{k}'s home band --- so verb-heavy clusters rise to the top of
their band without colliding with a neighbouring band. This preserves
both the domain orthogonality of Sub-Bands and the verb-slot phase-lock
property from the prior implementation.

\textbf{Consequence.} Two clusters in different bands share almost no
common Kuramoto phase drift, because their \texttt{ω} differs by at
least a band-width and the coupling terms
\texttt{sin(θ\_j\ −\ θ\_i)·A\_j} average to zero over any window longer
than the beat period. Domain interference during ordinary conversation
is suppressed by frequency orthogonality alone --- no relation table, no
cluster tag, no rule.

\paragraph{1.2 Voltage-Gated L6-A
Tunnels}\label{voltage-gated-l6-a-tunnels}

L6-A tunnels are the analogical bridges discovered by structural
resonance between cross-cluster pairs (Phase 1.3 online-L6A-discovery of
the CMD line of work). In VGHT, every \emph{inter-cluster} tunnel ---
and every tunnel discovered by L6-A regardless of cluster membership ---
is marked \emph{voltage-gated}. Its force is scaled at every tick by a
sigmoidal conductance:

\begin{verbatim}
C(A_a, A_b) = C_max · σ(κ · (A_a · A_b − θ_critical))
\end{verbatim}

where \texttt{A\_a\ =\ \textbar{}x\_a\textbar{}} and
\texttt{A\_b\ =\ \textbar{}x\_b\textbar{}} are the instantaneous
displacements of the two endpoint cavities, \texttt{σ} is the logistic
function, \texttt{κ} is a steepness parameter (default \texttt{12}), and
\texttt{θ\_critical} is the tunneling threshold (default \texttt{0.05}
in the field's natural units). \texttt{C\_max\ =\ 1}.

The force delivered to endpoints is scaled elementwise by \texttt{C}:

\begin{verbatim}
F_a ← F_a · C(A_a, A_b)
F_b ← F_b · C(A_a, A_b)
\end{verbatim}

\textbf{Consequence.} When either endpoint is quiescent,
\texttt{A\_a\ ·\ A\_b\ ≈\ 0}, \texttt{C\ ≈\ σ(−κθ\_critical)\ ≈\ 0}, and
the bridge presents effectively infinite resistance. The two clusters
exchange no energy. When both endpoints are simultaneously excited to a
level whose product exceeds \texttt{θ\_critical}, \texttt{C\ →\ 1}
sharply and the bridge conducts fully. This is the
resonant/quantum-tunneling analog explicitly identified in the
motivating architecture: the door opens only when the room on both sides
is loud enough to make the door irrelevant.

Because both endpoints must resonate at once, spurious cross-domain
leakage --- a stray phase alignment on one side of a bridge --- cannot
open the channel. Only genuine dual excitation, produced by an input
whose structural resonance actually spans two domains, does so. This is
the polymath spark.

\paragraph{1.3 Global Viscosity (GABA
analog)}\label{global-viscosity-gaba-analog}

The system tracks its total kinetic energy each tick:

\begin{verbatim}
E_total = ½ Σ_i m_i · v_i²
\end{verbatim}

An effective friction coefficient is derived as a saturating function of
\texttt{E\_total}:

\begin{verbatim}
μ_eff = μ_0 + k_avalanche · (1 − exp(−E_total / E_scale))
\end{verbatim}

with defaults \texttt{k\_avalanche\ =\ 0.10}, \texttt{E\_scale\ =\ 5.0}.
\texttt{μ\_eff} replaces the constant friction in the cavity
velocity-decay step and in the tunnel-damping term of that tick.

\textbf{Consequence.} Under normal, sparse activation \texttt{E\_total}
is small, \texttt{μ\_eff\ ≈\ μ\_0}, and the field evolves as before.
Under an avalanche --- many nodes firing simultaneously through spurious
phase alignment --- \texttt{E\_total} rises rapidly, \texttt{μ\_eff}
saturates toward \texttt{μ\_0\ +\ k\_avalanche}, and the excess kinetic
energy is dissipated in the same tick. Only the field components whose
amplitudes are large enough to overcome the elevated friction survive;
the noise floor is damped out. This is the physical analog of GABAergic
inhibitory tone.

Unlike the raw \texttt{exp(E\_total)} formulation, the bounded
\texttt{1\ −\ exp(−E\_total\ /\ E\_scale)} form cannot diverge,
guaranteeing numerical stability under any input, however extreme.

\begin{center}\rule{0.5\linewidth}{0.5pt}\end{center}

\subsubsection{2. Interaction of the Three
Mechanisms}\label{interaction-of-the-three-mechanisms}

The three mechanisms are not independent tuning knobs; they are coupled
by design.

\begin{itemize}
\tightlist
\item
  \textbf{Sub-Bands} determine which cavities \emph{can} phase-lock
  together in the absence of external excitation. Same-band cavities
  freely couple; cross-band cavities decouple by frequency alone.
\item
  \textbf{Voltage-gated bridges} re-establish cross-band coupling ---
  but only when both endpoints are already excited to a level that
  represents genuine resonance with the input. Without such excitation,
  the bridge is inert regardless of the graph's connectivity.
\item
  \textbf{Global viscosity} suppresses false excitation. When a spurious
  cascade begins to raise \texttt{E\_total}, μ\_eff rises and the
  cascade dies before it can meet the
  \texttt{A\_a\ ·\ A\_b\ \textgreater{}\ θ\_critical} condition needed
  to open any bridge.
\end{itemize}

The joint regime is: \emph{within-band coupling is free and continuous,
cross-band coupling is possible but expensive, and pathological
cross-band cascades are physically forbidden by dissipation.} Domain
isolation and polymath bridging cease to be tradeoffs.

\begin{center}\rule{0.5\linewidth}{0.5pt}\end{center}

\subsubsection{3. Physical
Interpretation}\label{physical-interpretation}

Each of the three mechanisms has a direct biological analog and an
independent literature suggesting the same functional role.

\begin{enumerate}
\def\labelenumi{\arabic{enumi}.}
\tightlist
\item
  \textbf{Frequency-band orthogonality} appears in cortical cognition as
  the observation that different cognitive processes (attention, working
  memory, motor planning) show distinguishable EEG band signatures (α,
  β, γ). Cortical microcircuits couple preferentially within-band, and
  cross-band coupling is a marker of deliberate binding (theta-gamma
  phase-amplitude coupling).
\item
  \textbf{Voltage-gated channels} are the mechanism by which NMDA
  receptors implement coincidence detection: their conductance is
  sigmoidal in the product of presynaptic glutamate release and
  postsynaptic depolarization, mirroring
  \texttt{C(A\_a,\ A\_b)\ =\ σ(A\_a\ ·\ A\_b\ −\ θ)} structurally.
\item
  \textbf{GABAergic inhibition} provides the global tone that keeps
  cortical excitability in the sparse-firing regime. Loss of GABAergic
  inhibition produces exactly the avalanche pattern VGHT's viscosity
  term suppresses: dense synchronous firing across unrelated cortical
  areas, the neural correlate of a seizure.
\end{enumerate}

VGHT does not claim to be a model of any of these biological mechanisms.
It claims that a substrate exhibiting the \emph{dynamical signatures} of
all three simultaneously --- with only three scalar parameters ---
provides the minimum architectural surface area on which
domain-orthogonal cognition with sparse analogical bridging is a stable
regime, rather than an accidental artifact of parameter tuning.

\begin{center}\rule{0.5\linewidth}{0.5pt}\end{center}

\subsubsection{4. Implementation
Contract}\label{implementation-contract}

\begin{itemize}
\tightlist
\item
  \textbf{No if/else branches} in the tick loop that inspect a relation
  label, a domain tag, or a cluster identifier. All three mechanisms are
  pure scalar arithmetic on existing per-cavity or per-tunnel fields.
\item
  \textbf{No new I/O boundary code} for domain metadata. Sub-Bands are
  assigned at bind time from cluster id, computed by the union-find that
  HNCA already runs. Voltage gating is a per-tunnel flag set once at
  bind for cross-cluster tunnels (and for any L6-A tunnel regardless of
  cluster membership). Viscosity is derived each tick from the same
  \texttt{total\_kinetic} accumulator already maintained for the RIFA IC
  pass.
\item
  \textbf{All defaults ON.} VGHT is not a feature flag. The three
  mechanisms are the substrate's normal dynamics.
\end{itemize}

\begin{center}\rule{0.5\linewidth}{0.5pt}\end{center}

\subsubsection{5. Parameters}\label{parameters-1}

\begin{longtable}[]{@{}lll@{}}
\toprule\noalign{}
Symbol & Default & Meaning \\
\midrule\noalign{}
\endhead
\bottomrule\noalign{}
\endlastfoot
\texttt{B} & 6 & Number of domain sub-bands \\
\texttt{ω\_lo}, \texttt{ω\_hi} & 8, 32 rad/s & Band ladder range \\
\texttt{κ} & 12 & Voltage-gate steepness \\
\texttt{θ\_critical} & 0.05 & Voltage-gate tunneling threshold \\
\texttt{k\_avalanche} & 0.10 & Peak added friction under avalanche \\
\texttt{E\_scale} & 5.0 & Kinetic-energy saturation scale \\
\end{longtable}

These are the only knobs. All other quantities are derived.

\begin{center}\rule{0.5\linewidth}{0.5pt}\end{center}

\subsubsection{6. Expected Empirical
Signature}\label{expected-empirical-signature}

A VGHT-enabled substrate is expected to show:

\begin{itemize}
\tightlist
\item
  \textbf{Single-domain prompts} produce single-domain phase-lock
  emissions with no cross-domain leakage into any slot, even after many
  turns of a single-domain conversation. The wavefront gate and the
  sub-band orthogonality are jointly responsible.
\item
  \textbf{Cross-domain prompts} --- where the input drives cavities in
  two different bands to simultaneous high amplitude --- open the
  voltage-gated bridges between those bands within a few ticks,
  producing a genuine polymath emission that spans domains.
\item
  \textbf{Injection of a single garbage or nonsense token} produces at
  worst its own auto-spawned label in output; the avalanche suppressor
  prevents its stray phase drift from recruiting unrelated cavities.
\item
  \textbf{Latency to cross-domain emission} is measurable: it should
  correspond to the time required for the voltage-gate sigmoid to
  saturate under paired excitation, i.e., a few tens of ticks. This
  latency is the phenomenological signature of the polymath spark and is
  measurable with the existing readout infrastructure.
\end{itemize}

\begin{center}\rule{0.5\linewidth}{0.5pt}\end{center}

\subsubsection{7. Relation to Prior CMD
Work}\label{relation-to-prior-cmd-work}

VGHT subsumes and refines three earlier CMD components without
discarding them.

\begin{itemize}
\tightlist
\item
  Phase 2.2 rubber-tunnel band-pass filters are a special case of
  voltage-gated tunnels with \texttt{θ\_critical\ =\ 0} (always
  conducting) and a Gaussian gate on \texttt{ω} alone. VGHT's sigmoidal
  amplitude gate is strictly more expressive.
\item
  Phase 3.1 multi-ontology substrates rely on manually loaded disjoint
  axiom files to achieve domain separation. VGHT achieves the same
  isolation from a single ontology by frequency orthogonality alone, and
  further allows analogical bridging without touching the ontology.
\item
  The cluster-energy depletion term in phase-locked emission remains a
  fallback for the readout stage. When VGHT is fully active, its
  contribution to output diversity is redundant with the sub-band
  separation.
\end{itemize}

\begin{center}\rule{0.5\linewidth}{0.5pt}\end{center}

\subsubsection{8. Failure Modes}\label{failure-modes-3}

The design has three known failure modes worth documenting explicitly.

\begin{enumerate}
\def\labelenumi{\arabic{enumi}.}
\tightlist
\item
  If the cluster count is less than \texttt{B}, some bands go unused,
  which is harmless but wastes spectral capacity.
\item
  If \texttt{θ\_critical} is set too high, cross-domain bridges never
  open even for well-posed polymath prompts. If set too low, the system
  reduces to the pre-VGHT shared-oscillator field.
\item
  If \texttt{k\_avalanche} is set too high, meaningful large-amplitude
  emissions are damped alongside noise; the system becomes
  under-expressive. Empirical tuning against a held-out prompt set is
  suggested.
\end{enumerate}

\begin{center}\rule{0.5\linewidth}{0.5pt}\end{center}

\subsubsection{9. Summary}\label{summary-2}

VGHT closes the domain-isolation / polymath-bridging tradeoff by mapping
domains to orthogonal frequency bands, making cross-domain bridges
voltage-gated by joint amplitude resonance, and adding a global
energy-dependent viscosity that dissipates avalanches before they can
spuriously open bridges. All three mechanisms are scalar physics with a
total of six parameters. The tick loop remains free of if/else on
relations, domain tags, or cluster identifiers. The result is a
substrate that is simultaneously domain-locked in default operation and
polymath-capable when its input carries structural resonance across
bands --- the physical realization of the intuition that a polymath is
not a system that mixes domains randomly, but one whose domains are
ordinarily silent to each other and speak only when the same underlying
pattern excites both.

\newpage

\subsection{Chaotic Strange-Attractor Trajectory (CSAT): Breaking
Ground-State Locking in Synthetic
Intelligence}\label{chaotic-strange-attractor-trajectory-csat-breaking-ground-state-locking-in-synthetic-intelligence}

\textbf{Abstract} Pure wave-based SI architectures often suffer from
dictionary-locking, where an input triggers the absolute closest
topological match and collapses into static definitions. Humans do not
converse by printing definitions; human thought is an open-ended
trajectory through a semantic manifold. We introduce \textbf{Chaotic
Strange-Attractor Trajectory (CSAT)}. By adding a low-dimensional
chaotic force (driven by a native Lorenz or Rössler system integrated
directly into the substrate's global viscosity), the system is prevented
from settling into static ground-state energy wells. This forces the
phase-sweep projection to continuously dance along a strange attractor,
creating human-like contextual drifting without losing structural
coherence.

\begin{verbatim}
                  [ Input Concept: linux ]
                             │
                             ▼
                   [ Static Definition Well ] (Old Way: Stays at 'nano')
                             │
                    (CSAT Viscosity Push) ───► [ Lorenz Chaos Attractor ]
                             │
                             ▼
               [ Dynamic Trajectory Path ]
          (Emits: nano ──► terminal ──► bash ──► kernel)
\end{verbatim}

\subsubsection{1. The Core Issue: Infinite Potential
Wells}\label{the-core-issue-infinite-potential-wells}

When \texttt{linux} is presented as input, the memory potential at the
\texttt{linux} and \texttt{nano} cavities is high enough that the entire
field's energy collapses into them and is trapped there. In physical
terms the energy landscape forms a \textbf{deep potential well}. This
does not happen in the human brain, because an undercurrent of noise ---
\textbf{metastable dynamics} --- is always present, and it prevents any
node from remaining fixed.

\subsubsection{2. Novel Architecture: CSAT
Mechanisms}\label{novel-architecture-csat-mechanisms}

\paragraph{2.1. The Substrate Lorenz
Attractor}\label{the-substrate-lorenz-attractor}

We integrate a three-line, low-level nonlinear chaotic system --- the
Lorenz equations --- directly into the global \texttt{tick()} loop of
\texttt{physics.hpp}: \[ dx = \sigma(y - x) \cdot dt \]
\[ dy = (x(\rho - z) - y) \cdot dt \]
\[ dz = \tau \cdot (xy - \beta z) \cdot dt \]

\emph{(Note: all three derivatives carry the time-scale factor \(\tau\);
see Numerical Constraints below.)}

\textbf{Numerical constraints --- required for a working
implementation.}

\begin{enumerate}
\def\labelenumi{\arabic{enumi}.}
\item
  \textbf{Time-scale factor \(\tau\).} Forward-Euler integration of the
  Lorenz system at \(\sigma = 10\), \(\rho = 28\), \(\beta = 8/3\) is
  unstable for \(dt \gtrsim 0.01\). The substrate runs at \(dt = 0.02\),
  past the stability boundary --- naive integration diverges to
  \texttt{NaN} within a few hundred ticks. The factor \(\tau\) (default
  \(0.3\)) slows the attractor without altering its geometry, restoring
  stability \emph{and} giving direct control over the drift timescale.
  \(\tau\) should be chosen so the attractor's orbital period is
  comparable to one conversational turn (\textasciitilde200 ticks). If
  the attractor cycles much faster, the chaotic push averages out within
  a single settle window and contributes nothing.
\item
  \textbf{Amplitude normalization.} The Lorenz \(x\) coordinate ranges
  over approximately \([-20, +20]\), not \([-1, +1]\). All downstream
  use must reference the normalized coordinate \(\hat{x} = x / 20\).
  Using raw \(x\) in a multiplicative viscosity term drives effective
  friction \emph{negative} for any gain above \(\gamma \approx 0.05\)
  --- anti-damping, unbounded energy injection, field explosion. This is
  a silent failure mode: no error is raised, only a wall of \texttt{NaN}
  several hundred ticks later.
\end{enumerate}

\begin{itemize}
\tightlist
\item
  There is no neural network and no training here. This is pure physical
  mathematics, updating three scalar variables (\(x, y, z\)) per tick.
\end{itemize}

\paragraph{2.2. Dynamic Viscosity Modulation (The
Push)}\label{dynamic-viscosity-modulation-the-push}

The VGHT layer carries a constant viscosity that suppresses avalanches.
We modulate that viscosity dynamically by the \(x\) coordinate of the
Lorenz attractor:
\[ \mu_{\text{dynamic}} = \mu_{\text{base}} \cdot (1.0 + \gamma \cdot \hat{x}_{\text{lorenz}}) \]

\textbf{However --- this global form provably cannot work.} See §2.2.1.
The corrected per-cavity form is given in §2.2.2.

\paragraph{2.2.1 Negative result: global viscosity modulation is
rank-preserving}\label{negative-result-global-viscosity-modulation-is-rank-preserving}

Applied as a \textbf{global} scalar in the cavity velocity-decay step
\(v_i \leftarrow v_i(1 - \mu_t)\), over \(N\) ticks:

\[ v_i(N) \;=\; v_i(0) \prod_{t=1}^{N}\bigl(1 - \mu_t\bigr) \;=\; v_i(0)\cdot C \]

where \(C\) depends only on the schedule \(\{\mu_t\}\) and is therefore
\textbf{identical for every cavity}. Hence

\[ \frac{v_i(N)}{v_j(N)} \;=\; \frac{v_i(0)}{v_j(0)} \qquad \forall\, i,j. \]

Amplitude \emph{ratios} are exactly invariant under any global viscosity
schedule, chaotic or otherwise. Because CPME §1.6 phase-collapse
emission selects by arg-max over a weight monotone in amplitude, the
emitted winner is unchanged. A globally-modulated viscosity alters how
\emph{loud} the field is, never \emph{which cavity is loudest}.

Concretely: if \texttt{think} carries 10× the amplitude of
\texttt{terminal} before the chaotic push, it carries 10× after.
\texttt{linux\ →\ nano} remains \texttt{linux\ →\ nano}. Multiplying the
entire field by a wiggling scalar is a monotone transform on the
ranking, and monotone transforms cannot reorder an arg-max.

This is a genuine negative result and it constrains the design:
\textbf{any mechanism intended to break ground-state locking must couple
to per-cavity state.}

\paragraph{2.2.2 The working form --- heat-coupled chaotic
viscosity}\label{the-working-form-heat-coupled-chaotic-viscosity}

The substrate already maintains a per-cavity scalar meaning exactly
``this cavity has been active too long'': the CMD thermodynamic
habituation heat \(T_i\). Coupling the chaotic term to \(T_i\) makes the
push selective:

\[ \mu_i \;=\; \underbrace{\mu_{\text{base}} + k_{\text{aval}}\bigl(1 - e^{-E_{\text{tot}}/E_s}\bigr)}_{\text{VGHT §1.3, unchanged}} \;+\; \underbrace{\gamma \cdot \hat{x}_{\text{lorenz}} \cdot \tanh\!\bigl(T_i / T_{\text{crit}}\bigr)}_{\text{CSAT}} \]

Entrenched cavities (high \(T_i\)) absorb the chaotic shake; fresh or
cold cavities barely register it. Rank can now flip, and that is genuine
metastability rather than uniform loudness modulation.

The \(\tanh\) bounds the term so \(\mu_i\) cannot run away. Because
\(\hat{x}\) alternates sign along the attractor, an entrenched cavity
experiences alternating extra-damping and reduced-damping --- it is
\emph{shaken}, not merely suppressed. Shaking permits escape;
suppression would only silence.

\textbf{Composition with the Curie gate.} \(T_i\) already gates tunnel
coupling via \(\exp(-T_i/T_{\text{crit}})\). Under CSAT the same state
variable acquires a second role: \emph{talking too much} now implies
both \emph{losing coupling} and \emph{becoming subject to chaotic
eviction}. One state variable, two independent escape mechanisms, no new
bookkeeping, tick loop still branchless.

\begin{itemize}
\tightlist
\item
  \textbf{Mechanism.} Whenever the system begins to lock onto a
  particular word (a ground state), the Lorenz chaotic push perturbs the
  global viscosity. This drives the field off that cavity and forces
  energy to spread into the cavities connected to it ---
  \texttt{terminal}, \texttt{bash}, \texttt{kernel}.
\end{itemize}

\paragraph{\texorpdfstring{2.3. Trajectory Momentum --- chaos supplies
\emph{timing}, CPSB supplies
\emph{direction}}{2.3. Trajectory Momentum --- chaos supplies timing, CPSB supplies direction}}\label{trajectory-momentum-chaos-supplies-timing-cpsb-supplies-direction}

\textbf{The division of labour must be stated precisely.} The Lorenz
system has no knowledge of the graph: \(x(t)\) does not know which
cavity is \texttt{terminal} and which is \texttt{bash}. The attractor's
geometry lives in its own \((x,y,z)\) space and bears no relation to the
semantic manifold. It therefore cannot, by itself, produce a
semantically coherent chain.

What CSAT contributes is \textbf{escape timing} --- \emph{when} to leave
the current basin. What produces the ordered chain
\texttt{linux\ →\ terminal\ →\ bash\ →\ kernel} is \textbf{CPSB gauge
coupling plus graph topology} --- \emph{where} the energy goes once it
leaves.

\begin{quote}
\textbf{CSAT provides escape timing. CPSB provides escape direction.}
\end{quote}

This is a clean, non-overlapping role for CSAT within the stack, and it
is the honest claim. Attributing semantic coherence to the attractor
itself would overstate the mechanism; the coherence is already earned by
CPSB and the ontology.

Chaos is not randomness. The Lorenz attractor traces a definite
geometric pattern --- a \emph{strange attractor}. The consequence is
that the signal does not emit random words: starting from \texttt{linux}
it builds a chain that respects the structural rules supplied by CPSB,
\texttt{linux} \(\rightarrow\) \texttt{terminal} \(\rightarrow\)
\texttt{bash} \(\rightarrow\) \texttt{kernel}.

\begin{center}\rule{0.5\linewidth}{0.5pt}\end{center}

\subsubsection{3. Parameters}\label{parameters-2}

\begin{longtable}[]{@{}
  >{\raggedright\arraybackslash}p{(\linewidth - 4\tabcolsep) * \real{0.3333}}
  >{\raggedright\arraybackslash}p{(\linewidth - 4\tabcolsep) * \real{0.3333}}
  >{\raggedright\arraybackslash}p{(\linewidth - 4\tabcolsep) * \real{0.3333}}@{}}
\toprule\noalign{}
\begin{minipage}[b]{\linewidth}\raggedright
Symbol
\end{minipage} & \begin{minipage}[b]{\linewidth}\raggedright
Default
\end{minipage} & \begin{minipage}[b]{\linewidth}\raggedright
Meaning
\end{minipage} \\
\midrule\noalign{}
\endhead
\bottomrule\noalign{}
\endlastfoot
\(\sigma, \rho, \beta\) & 10, 28, 8/3 & Standard Lorenz coefficients \\
\(\tau\) & 0.30 & Time-scale factor (stability + drift rate) \\
\(\gamma\) & 0.04 & Chaotic viscosity gain on \(\hat{x}\) \\
\(\hat{x}\) & \(x/20\) & Normalized Lorenz coordinate \\
\texttt{csat\_enabled} & see §4.1 & Master switch --- \textbf{not}
unconditionally on \\
\end{longtable}

Five parameters. Setting \(\gamma = 0\) recovers the pre-CSAT substrate
exactly.

\begin{center}\rule{0.5\linewidth}{0.5pt}\end{center}

\subsubsection{4. Failure Modes}\label{failure-modes-4}

\paragraph{4.1 Lyapunov divergence breaks
bit-reproducibility}\label{lyapunov-divergence-breaks-bit-reproducibility}

Chaos is deterministic, so seeded reproducibility survives \emph{in
principle}: the same \texttt{bind\_seed} and the same Lorenz initial
conditions yield the same trajectory. But the Lorenz system has a
positive Lyapunov exponent --- a \(10^{-7}\) floating-point difference
grows to \(O(1)\) within roughly a thousand ticks. Different compiler
flags, \texttt{-ffast-math}, or a different FPU will therefore produce
divergent trajectories from identical seeds.

This directly conflicts with the bit-exactness that the 20-seed RB+G
table (\texttt{rbg.md} §6.1) and the four-lever ablation grid (§6.2)
depend on.

\textbf{CSAT is consequently the first mechanism in the stack that must
not default unconditionally ON.} Recommended policy:

\begin{itemize}
\tightlist
\item
  \texttt{csat\_enabled\ =\ true} in the conversational shell
  (\texttt{si\_chat})
\item
  \texttt{csat\_enabled\ =\ false} forced inside
  \texttt{study\_recall\_seeds} and \texttt{study\_recall\_ablation}
\end{itemize}

Published numbers stay bit-reproducible; conversation still gets the
drift. This is a documented exception to the ``all defaults ON''
contract in the implementation guide §11.3, and the exception should be
recorded there.

\paragraph{4.2 CSAT does not address stale amplitude
accumulation}\label{csat-does-not-address-stale-amplitude-accumulation}

Ground-state locking has more than one cause, and CSAT addresses only
the genuine-potential-well case. In particular, with friction
\(\mu_0 = 0.006\) per tick, roughly \(30\%\) of a cavity's amplitude
survives a 200-tick turn (\(0.994^{200} \approx 0.30\)). A cavity driven
once at high amplitude therefore remains dominant for several subsequent
turns by pure arithmetic, with no potential well involved.

That is a \emph{dissipation} deficiency, not a chaos deficiency. Chaotic
shaking of an over-energised field does not remove the excess energy.
Fixing it requires either stronger inter-turn damping or an explicit
energy budget per turn --- orthogonal to CSAT and left as separate work.

\paragraph{4.3 Chaos on a sparse graph shakes an empty
room}\label{chaos-on-a-sparse-graph-shakes-an-empty-room}

If the driven concept has few or no outbound tunnels, there is no
trajectory for the attractor to push the field along. CSAT then produces
no observable change: the escape timing fires, but there is nowhere to
escape to. CSAT is therefore a mechanism whose benefit scales with
ontology density, and it should be evaluated only on subgraphs rich
enough to support multi-hop propagation.

\begin{center}\rule{0.5\linewidth}{0.5pt}\end{center}

\subsubsection{5. Relation to the Existing
Stack}\label{relation-to-the-existing-stack}

\begin{longtable}[]{@{}
  >{\raggedright\arraybackslash}p{(\linewidth - 4\tabcolsep) * \real{0.3333}}
  >{\raggedright\arraybackslash}p{(\linewidth - 4\tabcolsep) * \real{0.3333}}
  >{\raggedright\arraybackslash}p{(\linewidth - 4\tabcolsep) * \real{0.3333}}@{}}
\toprule\noalign{}
\begin{minipage}[b]{\linewidth}\raggedright
Layer
\end{minipage} & \begin{minipage}[b]{\linewidth}\raggedright
Mechanism
\end{minipage} & \begin{minipage}[b]{\linewidth}\raggedright
Role
\end{minipage} \\
\midrule\noalign{}
\endhead
\bottomrule\noalign{}
\endlastfoot
CMD & Curie gate on \(T_i\) & entrenched cavity loses \emph{coupling} \\
\textbf{CSAT} & \textbf{chaotic viscosity on \(T_i\)} &
\textbf{entrenched cavity loses \emph{stability}} \\
VGHT & sub-bands, voltage gates, viscosity & domain isolation +
avalanche control \\
CPME & momentum, acoustic mass, settle, sweep & emission eligibility and
ordering \\
CPSB & gauge \(\Psi\) on tunnels & escape \emph{direction} / word
order \\
CPME-II & phase-consistent Hebbian & stores the phase gradient across
turns \\
\end{longtable}

CSAT and the CMD Curie gate read the same state variable \(T_i\) but act
on different quantities --- coupling strength versus dynamical
stability. They are complementary, not redundant.

\begin{center}\rule{0.5\linewidth}{0.5pt}\end{center}

\newpage

\subsection{Morphological Resonance Synthesis (MRS): Language
Construction Without Tokens, Gradients, or Learned
Embeddings}\label{morphological-resonance-synthesis-mrs-language-construction-without-tokens-gradients-or-learned-embeddings}

\emph{The byte-lane sensory layer is currently input-only: bytes enter
the substrate through learned couplings, but output leaves it as
retrieved label strings. That retrieval is the last token-shaped
operation in the pipeline. MRS removes it by making the byte layer
bidirectional and giving it the same three phase mechanisms the concept
layer already has --- chirality on input, phase-consistent
consolidation, and phase-collapse projection on output. Word forms are
then} synthesised \emph{by sweeping a byte-phase field rather than}
retrieved \emph{from a table.}

\begin{center}\rule{0.5\linewidth}{0.5pt}\end{center}

\subsubsection{0. The Remaining
Asymmetry}\label{the-remaining-asymmetry}

The CMD substrate ingests raw bytes. There is no tokeniser: each of 256
\texttt{ByteLane} oscillators is driven by its byte value, and lanes
reach the concept field only through Hebbian couplings \(w[b][c]\) that
begin at zero. On the input side, the architecture is already
sub-lexical and token-free.

The output side is not. \texttt{readout\_phase\_collapse()} selects
concept cavities by phase and then prints
\texttt{si.graph.nodes{[}label\_key{]}.label} --- a \textbf{stored
string}. The physics decides \emph{which} concepts speak and \emph{in
what order}, but the \emph{form of each word} is a lookup.

This asymmetry has three consequences:

\begin{enumerate}
\def\labelenumi{\arabic{enumi}.}
\tightlist
\item
  \textbf{No morphology.} The engine can emit \texttt{walk} if
  \texttt{walk} is a stored label, and \texttt{walked} if
  \texttt{walked} is separately stored, but it cannot construct
  \texttt{walked} from \texttt{walk}. There is no mechanism by which
  stem and inflection compose.
\item
  \textbf{No novel forms.} Any word form absent from the label table is
  unreachable, regardless of how strongly the substrate's state implies
  it.
\item
  \textbf{A token-shaped residue.} The label table is a finite inventory
  of atomic surface strings indexed by graph key. Functionally that is a
  vocabulary, and its presence is what makes the pipeline token-like at
  the boundary even though the interior is not.
\end{enumerate}

MRS closes the asymmetry. The mechanism is a direct reuse --- one layer
down --- of machinery the substrate already contains.

\begin{center}\rule{0.5\linewidth}{0.5pt}\end{center}

\subsubsection{1. Architecture}\label{architecture-2}

Let the byte-lane population be \(L_b\) for \(b \in [0,256)\), each an
oscillator with displacement \(x_b\), velocity \(v_b\), and natural
frequency \(\omega_b = 3.0 + b/256\) (as in the existing CMD sensory
layer). Let \(w[b][c]\) be the learned byte→concept binding weights.

MRS adds four mechanisms.

\paragraph{1.1 Reverse Drive (concept →
lane)}\label{reverse-drive-concept-lane}

The existing forward path applies a force from lane to concept:

\[ F_c \mathrel{+}= \frac{w[b][c] \cdot x_b}{m_c} \]

MRS adds the transpose, using the \textbf{same weights}:

\[ F_b \mathrel{+}= \eta_{\text{rev}} \cdot w[b][c] \cdot x_c \]

with \(\eta_{\text{rev}}\) a reverse-gain scalar (default \(0.4\)). No
new storage: the coupling matrix is used in both directions, making the
byte↔concept layer a bidirectional associative memory implemented as
coupled oscillators rather than as a matrix multiplication.

Consequence: when concept cavities resonate, the byte lanes whose
bindings to those concepts are strongest resonate in turn. The lanes
that light up are the letters of the words those concepts were
historically written with.

\paragraph{1.2 Lane Chirality}\label{lane-chirality}

Byte lanes acquire the CPME §1.1 arrival timestamp:

\[ t_{\text{arrival},b} \leftarrow t_{\text{sim}} \quad \text{on each drive of } L_b \]

During input of \texttt{hi}, lane \(L_{\texttt{h}}\) is driven before
\(L_{\texttt{i}}\) and therefore accumulates more phase. Letter order
within a word becomes a phase gradient across the lane population ---
exactly the mechanism CPME already uses for word order across concepts,
applied one level down.

\paragraph{1.3 Lane-to-Lane Gauge
Tunnels}\label{lane-to-lane-gauge-tunnels}

A sparse set of lane↔lane tunnels is created on co-occurrence, each
carrying a CPSB-style gauge offset \(\Psi_{bb'}\) and consolidated by
the CPME-II phase-consistent rule:

\[ \Delta w_{bb'} \;\propto\; v_b \cdot v_{b'} \cdot \cos\!\bigl(\theta_b - \theta_{b'} - \Psi_{bb'}\bigr) \]

where \(\theta_b = \operatorname{atan2}(v_b, \omega_0 x_b)\). Unlike the
concept layer, \(\Psi_{bb'}\) is \textbf{not} assigned from a relation
table --- it is \emph{learned} as the running circular mean of observed
phase differences:

\[ \Psi_{bb'} \leftarrow \arg\Bigl( \lambda\, e^{i\Psi_{bb'}} + (1-\lambda)\, e^{i(\theta_b - \theta_{b'})} \Bigr) \]

This is the substrate's representation of \emph{sequence}: the tunnel
between \texttt{h} and \texttt{i} stores the typical phase lead of
\texttt{h} over \texttt{i}. Nothing about letters, words, or grammar is
encoded declaratively; only phase offsets between oscillators.

\paragraph{1.4 Byte-Level Phase-Collapse
Projection}\label{byte-level-phase-collapse-projection}

Output is generated by the CPME §1.6 sweep applied to the lane
population. Carrier phase \(\phi_{\text{carrier}}\) advances
monotonically across \([0, 2\pi)\) in \(K_{\text{byte}}\) steps; at each
step the emitted byte is the arg-max of

\[ w_b(\phi_{\text{carrier}}) \;=\; |x_b| \cdot \operatorname{ReLU}(P_b) \cdot \exp\!\bigl(-\kappa_b (\phi_{\text{carrier}} - \theta_b)^2\bigr) \]

with \(P_b\) the lane's envelope momentum as in CPME §1.2.

\textbf{Critical divergence from the concept sweep.} In the concept
layer, a winner is suppressed from later sweep steps --- a concept
appears once in a clause. In the byte layer, \textbf{winners are not
suppressed}, and the sweep may make multiple passes. Letters repeat
within words (\texttt{banana}, \texttt{letter}, \texttt{success}); a
suppression rule would make repeated graphemes unreachable. The
difference is physically motivated by what each layer represents, not
introduced as an exception.

Word boundaries emerge from amplitude gaps: when the swept weight falls
below the emission floor for a contiguous span of carrier phase, a
boundary is emitted. Boundaries are therefore not delimiters in a
grammar but silences in a resonance field.

\begin{center}\rule{0.5\linewidth}{0.5pt}\end{center}

\subsubsection{2. Why This Is Not Tokens, ML, DL, or an
LLM}\label{why-this-is-not-tokens-ml-dl-or-an-llm}

\textbf{Not tokenisation.} There is no vocabulary, no merge table, no
subword inventory, no \texttt{\textless{}unk\textgreater{}}. Input is
256 byte lanes; output is a swept byte sequence. The only inventory is
the byte alphabet itself, which is a property of the encoding, not a
learned artefact.

\textbf{Not machine learning in the gradient sense.} There is no loss
function, no gradient, no backpropagation, no optimiser, no train/test
split. All adaptation is the local Hebbian rule of §1.3, which uses only
quantities available at the two endpoints of a single tunnel at a single
tick. Nothing computes a derivative of a global objective.

\textbf{Not deep learning.} There is no layer stack in the DL sense ---
no weight matrices applied in sequence, no nonlinearity between layers,
no depth over which credit is assigned. The byte layer and concept layer
are two populations of coupled oscillators exchanging forces
bidirectionally through one shared coupling matrix.

\textbf{Not a language model.} No distribution over next symbols is ever
formed. There is no softmax, no sampling temperature, no autoregressive
conditioning on prior output. The emitted sequence is the projection of
a static field state onto a swept phase axis --- a geometric readout,
not a sequential prediction. Given the same field state the same
sequence is produced deterministically, and the sequence is generated
\emph{simultaneously} in the sense that all lanes are already in their
phases before the sweep begins.

\textbf{What it is:} a bidirectional associative resonance memory
between a 256-oscillator sub-lexical layer and an \(N\)-oscillator
conceptual layer, read out by phase-swept projection.

\begin{center}\rule{0.5\linewidth}{0.5pt}\end{center}

\subsubsection{3. Where the Storage Lives (Translation Capacity
Lemma)}\label{where-the-storage-lives-translation-capacity-lemma}

The Translation Capacity Lemma (SI §3.1) states that any system
discriminating \(N\) outcomes must store \(\log_2 N\) bits somewhere.
MRS does not evade this and does not claim to. The storage relocates:

\begin{longtable}[]{@{}
  >{\raggedright\arraybackslash}p{(\linewidth - 2\tabcolsep) * \real{0.5000}}
  >{\raggedright\arraybackslash}p{(\linewidth - 2\tabcolsep) * \real{0.5000}}@{}}
\toprule\noalign{}
\begin{minipage}[b]{\linewidth}\raggedright
Before MRS
\end{minipage} & \begin{minipage}[b]{\linewidth}\raggedright
After MRS
\end{minipage} \\
\midrule\noalign{}
\endhead
\bottomrule\noalign{}
\endlastfoot
\(N\) label strings in \texttt{graph.nodes{[}{]}.label} &
\(256 \times N\) byte→concept couplings \(w[b][c]\) \\
--- & \(\le 256^2\) lane↔lane couplings and offsets \(\Psi_{bb'}\) \\
\end{longtable}

The relocation is the substantive change. A label table stores each
surface form \emph{atomically and independently}: \texttt{walk} and
\texttt{walked} share no representation, so knowing one contributes
nothing to producing the other. A coupling matrix stores forms
\emph{distributedly and compositionally}: the lanes and phase offsets
that produce \texttt{walk} are largely the same ones that produce
\texttt{walked}, differing by the additional offsets contributed by the
co-active inflectional concept.

Generalisation to unseen forms therefore becomes possible in principle
--- not because information was created, but because the information was
stored in a factored representation instead of an atomic one. The bit
count is not reduced; its \emph{organisation} is what changes.

\begin{center}\rule{0.5\linewidth}{0.5pt}\end{center}

\subsubsection{4. Expected Dynamics}\label{expected-dynamics-4}

\textbf{Prompt:} \texttt{metformin\ lowers\ glucose}

\begin{enumerate}
\def\labelenumi{\arabic{enumi}.}
\tightlist
\item
  \textbf{Byte injection.} Lanes for each character are driven in
  arrival order; \(t_{\text{arrival}}\) records the chirality gradient.
  Lane↔lane tunnels consolidate the observed offsets (\texttt{m} leads
  \texttt{e} leads \texttt{t}\ldots).
\item
  \textbf{Concept resonance.} Forward couplings \(w[b][c]\) excite the
  \texttt{metformin}, \texttt{lowers}, \texttt{glucose} cavities;
  propagation through the ontology raises downstream cavities
  (\texttt{ampk}, \texttt{glucose\_uptake}).
\item
  \textbf{Acoustic-mass muting.} Directly-driven concepts are muted
  (CPME §1.3); downstream concepts are not.
\item
  \textbf{Reverse drive.} The surviving downstream concepts drive byte
  lanes through \(w[b][c]^{\!\top}\). Lanes corresponding to the letters
  of \texttt{ampk}, \texttt{glucose}, \texttt{uptake} resonate ---
  \emph{without those strings being retrieved}.
\item
  \textbf{Byte sweep.} The carrier phase sweeps the lane population.
  Lanes emit in phase order, with boundaries at amplitude gaps. Output
  is a byte sequence assembled from resonance.
\end{enumerate}

\textbf{Morphological composition.} If the inflectional concept for
past-tense has been bound to the lanes \texttt{e},\texttt{d} through
prior exposure, and it is co-active with \texttt{lower}, the reverse
drive raises both the stem lanes and the suffix lanes. The sweep ---
ordered by learned phase offsets --- emits \texttt{lowered}, a form that
need never have been stored as a label.

\begin{center}\rule{0.5\linewidth}{0.5pt}\end{center}

\subsubsection{4b. Empirical Result: Single-Byte Lanes Are Insufficient
(measured)}\label{b.-empirical-result-single-byte-lanes-are-insufficient-measured}

The single-byte formulation of §1.1--§1.4 was implemented and tested.
\textbf{It does not work, and the reason is representational rather than
a matter of tuning.}

Recognition was attempted by grounding every concept label to the 256
lanes at load time (binding weight = inverse byte document-frequency
across the lexicon) and then recognising a query by driving its bytes
and reading which concepts accumulated energy. Measured on a 2,240
-concept graph:

\begin{longtable}[]{@{}lll@{}}
\toprule\noalign{}
query & recognised as & correct \\
\midrule\noalign{}
\endhead
\bottomrule\noalign{}
\endlastfoot
\texttt{metfromin} & \texttt{from}, \texttt{front}, \texttt{rfe} & ✗ \\
\texttt{metformin} & \texttt{from}, \texttt{front}, \texttt{rfe} & ✗ \\
\texttt{insluin} & \texttt{uni}, \texttt{ull}, \texttt{sil} & ✗ \\
\texttt{fissoin} & \texttt{fission} & ✓ \\
\end{longtable}

Two distinct causes, and separating them matters:

\textbf{(a) A normalisation error.} Binding weight was divided by
\(\sqrt{|\text{label}|}\). For the query \texttt{metformin}, the
substring \texttt{from} scored \(S/2\) while \texttt{metformin} itself
scored \(S/3\) --- the extra bytes \(\{e,t,i,n\}\) are common and
contribute ≈ 0 weight, but the length penalty applies regardless. A
substring therefore outscored the whole word, which is why recognition
failed even on \emph{correctly spelled} input. This was a genuine bug.

\textbf{(b) A representational limit.} Removing the length division and
replacing it with symmetric coverage (harmonic mean of query-recall and
label-precision over byte sets) fixed the substring case and broke a
different one: long labels sharing many bytes with the query
(\texttt{intent\_from\_marker\_label}) began to win, because a concept
bound to more lanes accumulates more raw energy regardless of coverage.
Each subsequent scoring adjustment relocated the failure rather than
removing it.

That pattern --- a defect that moves under every reweighting but never
disappears --- is the signature of insufficient representation. The
decisive case is \texttt{metformin} versus \texttt{metfromin}:
\textbf{they have identical byte sets.} A bag of distinct bytes cannot
distinguish a word from its own transposition, so no weighting function
over that bag can succeed. This is §5.3 (homograph collision)
manifesting not as an edge case but as the dominant failure.

\paragraph{4b.1 Required correction: bigram
lanes}\label{b.1-required-correction-bigram-lanes}

The sensory layer must carry ordering information. Replacing the 256
single-byte lanes with \(256^2\) ordered-pair lanes supplies it:

\[
\texttt{metformin} \rightarrow \{\texttt{me},\texttt{et},\texttt{tf},\texttt{fo},\texttt{or},\texttt{rm},\texttt{mi},\texttt{in}\}
\] \[
\texttt{metfromin} \rightarrow \{\texttt{me},\texttt{et},\texttt{tf},\texttt{fr},\texttt{ro},\texttt{om},\texttt{mi},\texttt{in}\}
\]

Five of eight pairs are shared; three differ. The transposition is now
\emph{visible to the substrate}, and graded recognition becomes
possible: a near-miss resonates strongly but not identically, which is
the behaviour §1.4 requires and the single-byte layer cannot deliver.

\textbf{This changes nothing about the architecture's commitments.} A
bigram lane is one oscillator per ordered character pair, with
displacement, velocity and natural frequency --- the same
\texttt{ByteLane} structure, a wider array. Specifically it introduces:

\begin{itemize}
\tightlist
\item
  no loss function, no gradient, no fitted parameter;
\item
  no attention, no softmax, no learned projection;
\item
  no probability semantics --- lane energy is displacement in a field.
\end{itemize}

The distinction that governs this architecture is not \emph{``are there
numbers''} but \emph{``where do the numbers come from.''} A binding
weight here is \textbf{computed} from a counted structural property of
the lexicon by a stated formula; changing the lexicon changes it
deterministically. A weight in a trained network is \textbf{fitted} by
optimising a loss and admits no explanation beyond that optimisation.
Descriptive measurement is not fitted inference --- measuring a mass
does not make mechanics machine learning.

n-gram indexing is also not novel-to-ML prior art: it originates in
1970s information retrieval (approximate string matching, spell
correction) and predates statistical learning entirely.

\paragraph{4b.2 Measured result after implementing bigram
lanes}\label{b.2-measured-result-after-implementing-bigram-lanes}

Bigram lanes were implemented (\texttt{PhysicsEngine::bigram\_bind\_w},
\texttt{bind\_bigram}). Grounding is identical in form to the
single-byte version --- inverse lane document-frequency across the
lexicon --- but over ordered pairs bracketed by BOS/EOS sentinels, so
word-initial and word-final position are themselves observable pairs.

Recognition divides energy received by a concept by the geometric mean
of its own binding mass and the query's:

\[
\text{resonance}(c) \;=\; \frac{\sum_{g \in Q \cap B(c)} w_g}{\sqrt{M_c \cdot M_Q}}
\]

This is energy per unit coupling --- the same reason acceleration is
force divided by mass --- and it is why neither long nor short labels
win by size. The earlier failures came from omitting exactly this
normalisation.

Measured on the same 2,240-concept graph:

\begin{longtable}[]{@{}lllll@{}}
\toprule\noalign{}
query & perturbation & recognised as & score & correct \\
\midrule\noalign{}
\endhead
\bottomrule\noalign{}
\endlastfoot
\texttt{metformin} & none & \texttt{metformin} & 1.000 & ✓ \\
\texttt{metfromin} & transposition & \texttt{metformin} & 0.686 & ✓ \\
\texttt{insluin} & transposition & \texttt{insulin} & 0.471 & ✓ \\
\texttt{fissoin} & transposition & \texttt{fission} & 0.621 & ✓ \\
\texttt{fisssion} & insertion & \texttt{fission} & 1.000 & ✓ \\
\texttt{cavty} & deletion & \texttt{cavity} & 0.830 & ✓ \\
\texttt{tunel} & deletion & \texttt{tunnel} & 0.885 & ✓ \\
\end{longtable}

\textbf{7 / 7}, against 1 / 4 for single-byte lanes.

The \texttt{metformin} / \texttt{metfromin} pair is the decisive one:
those two strings have \textbf{identical byte sets}, so the single-byte
formulation could not distinguish them under any weighting. Over bigrams
they share 5 of 8 pairs and differ in 3, and the correct concept is
recovered at 0.686 against 1.000 for the exact spelling.

That gradation is the property §1.4 requires and a hash lookup cannot
have: recognition degrades smoothly with distortion rather than failing
discontinuously. It is also the evidence that recognition is not string
matching --- no comparison between the query and any stored label occurs
at any point.

\textbf{Status:} bigram recognition is ENABLED in \texttt{codeq}, with
an acceptance floor of 0.35 so unknown tokens resonate with nothing
rather than matching noise. Single-byte recognition is retained in the
source only as the documented negative result. Per-query telemetry
reports \texttt{lexical} versus \texttt{resonant} recognition counts, so
residual dependence on exact lookup is a measured quantity rather than
an assumption.

\begin{center}\rule{0.5\linewidth}{0.5pt}\end{center}

\subsubsection{5. Failure Modes}\label{failure-modes-5}

\begin{enumerate}
\def\labelenumi{\arabic{enumi}.}
\item
  \textbf{Babble before consolidation.} With near-zero couplings the
  lane phases are unstructured and the sweep emits non-words. This is
  the correct initial state of a system that has not yet learned
  orthography, but it will read as garbage output until the lane↔lane
  offsets consolidate. Consolidation requires substantial exposure ---
  Hebbian rather than gradient, but still genuinely a training process
  measured in thousands of turns, not tens.
\item
  \textbf{Phase resolution limits word length.} With \(K_{\text{byte}}\)
  sweep steps the angular resolution is \(2\pi/K_{\text{byte}}\). A
  ten-character word requires ten sufficiently separated lane phases.
  Long words will be truncated or slurred until \(K_{\text{byte}}\) is
  raised, at proportional sweep cost.
\item
  \textbf{Homograph collision.} A single lane per byte value means the
  letter \texttt{a} has one phase at any instant. In a word requiring
  \texttt{a} at two different positions, the multi-pass sweep of §1.4
  recovers the repetition only if the two occurrences are separated by
  other letters whose phases bracket them. Words like \texttt{aa} or
  vowel-adjacent repeats are expected to fail. A positional lane bank
  (\(256 \times K\)) would resolve this at a large state cost and would
  reintroduce something functionally close to positional encoding; it is
  deliberately not proposed here.
\item
  \textbf{No orthographic well-formedness guarantee.} Nothing in the
  mechanism prevents emission of byte sequences that are not words.
  Well-formedness is a statistical consequence of consolidation, not a
  constraint. This is a genuine difference from a grammar-based
  generator and should be reported as such rather than patched with a
  dictionary filter --- a filter would reintroduce the vocabulary that
  MRS exists to remove.
\item
  \textbf{Reverse-gain instability.} Bidirectional coupling through a
  shared matrix creates a feedback loop: concepts drive lanes drive
  concepts. For \(\eta_{\text{rev}}\) above a critical value the loop
  has gain greater than unity and the field diverges. The VGHT §1.3
  avalanche viscosity bounds this, but \(\eta_{\text{rev}}\) must be
  tuned conservatively and the loop gain monitored.
\end{enumerate}

\begin{center}\rule{0.5\linewidth}{0.5pt}\end{center}

\subsubsection{6. Relation to the Existing
Stack}\label{relation-to-the-existing-stack-1}

MRS is the mirror of CPME one layer down. Every mechanism it uses
already exists in the substrate for the concept layer:

\begin{longtable}[]{@{}
  >{\raggedright\arraybackslash}p{(\linewidth - 4\tabcolsep) * \real{0.3333}}
  >{\raggedright\arraybackslash}p{(\linewidth - 4\tabcolsep) * \real{0.3333}}
  >{\raggedright\arraybackslash}p{(\linewidth - 4\tabcolsep) * \real{0.3333}}@{}}
\toprule\noalign{}
\begin{minipage}[b]{\linewidth}\raggedright
Mechanism
\end{minipage} & \begin{minipage}[b]{\linewidth}\raggedright
Concept layer
\end{minipage} & \begin{minipage}[b]{\linewidth}\raggedright
Byte layer (MRS)
\end{minipage} \\
\midrule\noalign{}
\endhead
\bottomrule\noalign{}
\endlastfoot
Chirality on input & \texttt{Cavity::t\_arrival} &
\texttt{ByteLane::t\_arrival} \\
Momentum filter & \texttt{env\_ema} on cavities & \texttt{env\_ema} on
lanes \\
Phase definition & \(\operatorname{atan2}(v_c, \omega_0 x_c)\) &
\(\operatorname{atan2}(v_b, \omega_0 x_b)\) \\
Gauge offset & \(\Psi\) from \texttt{psi\_for(rel)} & \(\Psi\)
\emph{learned} from observed phase \\
Phase-consistent Hebbian & CPME-II on tunnels & CPME-II on lane
tunnels \\
Phase-collapse sweep & emits concept labels & emits \textbf{bytes} \\
Winner suppression & yes (concept appears once) & \textbf{no} (letters
repeat) \\
\end{longtable}

The single genuinely new element is §1.1 reverse drive. Everything else
is the same physics applied to a different oscillator population ---
which is the strongest argument that the architecture is coherent rather
than assembled.

\begin{center}\rule{0.5\linewidth}{0.5pt}\end{center}

\subsubsection{7. Summary}\label{summary-3}

The CMD substrate was token-free on input and table-driven on output.
MRS removes the remaining table by making the byte-lane layer
bidirectional and equipping it with the chirality, consolidation, and
phase-collapse mechanisms already built for concepts. Word forms become
the projection of a byte-phase field rather than entries retrieved by
graph key.

No tokeniser, no gradient, no layer stack, no next-symbol distribution.
Storage obeys the Translation Capacity Lemma but is relocated from
atomic surface strings into a factored coupling matrix, which is what
makes morphological composition and novel forms reachable at all.

The honest cost is that consolidation is slow and early output is
babble. A system that constructs its own orthography must first fail to.

\end{document}
